\documentclass[11pt]{article}
\usepackage[margin=1in]{geometry}
\usepackage{amsmath,amssymb,amsthm,mathtools}
\usepackage{booktabs,longtable,array,graphicx}
\usepackage{xurl}
\usepackage{hyperref}
\usepackage{xcolor}
\usepackage{microtype}
\usepackage{enumitem}
\usepackage[T1]{fontenc}
\usepackage{lmodern}
\setlength{\emergencystretch}{2em}
\setcounter{tocdepth}{1}
\numberwithin{equation}{section}
\raggedbottom
\hypersetup{colorlinks=true,linkcolor=blue,citecolor=blue,urlcolor=blue}
\newtheorem{theorem}{Theorem}[section]
\newtheorem{proposition}[theorem]{Proposition}
\newtheorem{lemma}[theorem]{Lemma}
\newtheorem{corollary}[theorem]{Corollary}
\theoremstyle{remark}
\newtheorem{remark}[theorem]{Remark}
\theoremstyle{definition}
\newtheorem{question}[theorem]{Question}
\newtheorem*{executionhypothesis}{Packet execution hypothesis}
\theoremstyle{plain}
\newtheorem{theoremA}[theorem]{Theorem}
\newtheorem{theoremB}[theorem]{Theorem}
\newcommand{\OL}{\mathcal O_L}
\newcommand{\OM}{\mathcal O_M}
\newcommand{\Ql}{\mathbb Q}
\newcommand{\Z}{\mathbb Z}
\newcommand{\R}{\mathbb R}
\newcommand{\C}{\mathbb C}
\newcommand{\HH}{\mathbb H}
\newcommand{\Gal}{\operatorname{Gal}}
\newcommand{\End}{\operatorname{End}}
\newcommand{\Tr}{\operatorname{Tr}}
\newcommand{\Nm}{\operatorname{N}}
\newcommand{\GL}{\operatorname{GL}}
\newcommand{\SL}{\operatorname{SL}}
\newcommand{\extD}{\ensuremath{D_{\mathrm{CM}}}}
\newcommand{\extAlg}{\operatorname{Alg}}
\newcommand{\extHdg}{\operatorname{Hdg}}
\newcommand{\extHg}{\operatorname{Hg}}
\newcommand{\extsupp}{\operatorname{supp}}
\newcommand{\digest}[2]{\scalebox{0.9}{\texttt{#1#2}}}

\makeatletter
\renewcommand{\l@section}{\@dottedtocline{1}{0em}{2.3em}}
\makeatother
\newcommand{\externalref}[2]{\textnormal{\csname doccode@#1\endcsname.\ifcsname xnumber@#2\endcsname\csname xnumber@#2\endcsname\else ?\fi}}
% Generated from native auxiliary files; do not invent result numbers.
\expandafter\def\csname xnumber@app:arithdetail\endcsname{2}
\expandafter\def\csname xnumber@app:components\endcsname{7}
\expandafter\def\csname xnumber@app:computer\endcsname{13}
\expandafter\def\csname xnumber@app:orders\endcsname{8}
\expandafter\def\csname xnumber@app:prime43\endcsname{3}
\expandafter\def\csname xnumber@cor:12collapse\endcsname{11.5}
\expandafter\def\csname xnumber@cor:Gtransport\endcsname{5.5}
\expandafter\def\csname xnumber@cor:HCspecific\endcsname{5.12}
\expandafter\def\csname xnumber@cor:Weilalg\endcsname{5.9}
\expandafter\def\csname xnumber@cor:d7geometrycomplete\endcsname{9.12}
\expandafter\def\csname xnumber@cor:d7r83all\endcsname{9.9}
\expandafter\def\csname xnumber@cor:d7r83nontrivial\endcsname{9.7}
\expandafter\def\csname xnumber@cor:finitegeometric\endcsname{7.3}
\expandafter\def\csname xnumber@cor:geometry119\endcsname{9.15}
\expandafter\def\csname xnumber@cor:maximalgeometrycomplete\endcsname{11.10}
\expandafter\def\csname xnumber@cor:noAstablepp\endcsname{10.4}
\expandafter\def\csname xnumber@cor:noBstablepp\endcsname{4.4}
\expandafter\def\csname xnumber@cor:noCstablepp\endcsname{10.2}
\expandafter\def\csname xnumber@cor:orientedcollapse\endcsname{11.4}
\expandafter\def\csname xnumber@cor:rosatipi\endcsname{12.6}
\expandafter\def\csname xnumber@eight:cor:mt-splitting\endcsname{4.3}
\expandafter\def\csname xnumber@eight:cor:subfields\endcsname{4.2}
\expandafter\def\csname xnumber@eight:sec:lattice\endcsname{4.1}
\expandafter\def\csname xnumber@eight:sec:main\endcsname{4}
\expandafter\def\csname xnumber@eight:sec:sources\endcsname{4.2}
\expandafter\def\csname xnumber@eight:thm:all\endcsname{4.1}
\expandafter\def\csname xnumber@eq:42test\endcsname{6.4}
\expandafter\def\csname xnumber@eq:CMcriterion\endcsname{5.3}
\expandafter\def\csname xnumber@eq:Deltak\endcsname{8.1}
\expandafter\def\csname xnumber@eq:E1\endcsname{4.1}
\expandafter\def\csname xnumber@eq:HermitianE1\endcsname{5.1}
\expandafter\def\csname xnumber@eq:absnormalpha\endcsname{3.2}
\expandafter\def\csname xnumber@eq:adjdet\endcsname{12.1}
\expandafter\def\csname xnumber@eq:alpha0\endcsname{3.1}
\expandafter\def\csname xnumber@eq:companion\endcsname{6.1}
\expandafter\def\csname xnumber@eq:correct-charpoly\endcsname{12.2}
\expandafter\def\csname xnumber@eq:correct-equiv\endcsname{12.6}
\expandafter\def\csname xnumber@eq:correctschwarzcomposition\endcsname{12.9}
\expandafter\def\csname xnumber@eq:deltagamma\endcsname{12.5}
\expandafter\def\csname xnumber@eq:explicitcovergroup\endcsname{7.1}
\expandafter\def\csname xnumber@eq:galois-exponents\endcsname{12.7}
\expandafter\def\csname xnumber@eq:hgparamformula\endcsname{12.8}
\expandafter\def\csname xnumber@eq:mcmullenF\endcsname{6.3}
\expandafter\def\csname xnumber@eq:minpoly\endcsname{2.1}
\expandafter\def\csname xnumber@eq:normfiber\endcsname{4.1}
\expandafter\def\csname xnumber@eq:normpi\endcsname{12.3}
\expandafter\def\csname xnumber@eq:orbit-test\endcsname{6.5}
\expandafter\def\csname xnumber@eq:pi\endcsname{3.5}
\expandafter\def\csname xnumber@eq:principalpacketpoint\endcsname{3.1}
\expandafter\def\csname xnumber@eq:unitnormquotient\endcsname{11.1}
\expandafter\def\csname xnumber@eq:unitpi\endcsname{12.4}
\expandafter\def\csname xnumber@eq:vertexdisc\endcsname{5.2}
\expandafter\def\csname xnumber@eq:wi\endcsname{6.2}
\expandafter\def\csname xnumber@ext:eq:basis\endcsname{{$\mathcal  R$}}
\expandafter\def\csname xnumber@ext:lem:contraction\endcsname{2.4}
\expandafter\def\csname xnumber@ext:lem:lattice\endcsname{2.5}
\expandafter\def\csname xnumber@ext:sec:evidence\endcsname{2.6}
\expandafter\def\csname xnumber@ext:sec:inputs\endcsname{2.2}
\expandafter\def\csname xnumber@ext:sec:lattice\endcsname{2.4}
\expandafter\def\csname xnumber@ext:sec:main\endcsname{2}
\expandafter\def\csname xnumber@ext:thm:individual\endcsname{2.3}
\expandafter\def\csname xnumber@ext:thm:mixed\endcsname{2.1}
\expandafter\def\csname xnumber@hyp:packetexecution\endcsname{1.2}
\expandafter\def\csname xnumber@lem:adjugate\endcsname{12.5}
\expandafter\def\csname xnumber@lem:degree\endcsname{12.2}
\expandafter\def\csname xnumber@lem:narrowclass\endcsname{11.2}
\expandafter\def\csname xnumber@lem:packetpointfibers\endcsname{3.1}
\expandafter\def\csname xnumber@lem:primitivebinary\endcsname{5.10}
\expandafter\def\csname xnumber@lem:twochambers\endcsname{3.5}
\expandafter\def\csname xnumber@pp:cor:fermatprojector\endcsname{2.3}
\expandafter\def\csname xnumber@pp:cor:jacobianlist\endcsname{12.2}
\expandafter\def\csname xnumber@pp:cor:obstructedmodulifields\endcsname{6.3}
\expandafter\def\csname xnumber@pp:cor:quotientrealforms\endcsname{7.2}
\expandafter\def\csname xnumber@pp:cor:quotientweilhodge\endcsname{10.1}
\expandafter\def\csname xnumber@pp:cor:smallerjacobian\endcsname{13.2}
\expandafter\def\csname xnumber@pp:cor:underlyinggalois\endcsname{4.4}
\expandafter\def\csname xnumber@pp:cor:weilimage\endcsname{8.7}
\expandafter\def\csname xnumber@pp:eq:basis\endcsname{7.1}
\expandafter\def\csname xnumber@pp:eq:curve\endcsname{7.3}
\expandafter\def\csname xnumber@pp:eq:group\endcsname{7.2}
\expandafter\def\csname xnumber@pp:eq:unpolarizedlocaldata\endcsname{7.4}
\expandafter\def\csname xnumber@pp:prop:actualpolarizations\endcsname{7.12}
\expandafter\def\csname xnumber@pp:prop:allunderlyinggalois\endcsname{4.5}
\expandafter\def\csname xnumber@pp:prop:balancedtheta\endcsname{13.1}
\expandafter\def\csname xnumber@pp:prop:bcisogeny\endcsname{7.5}
\expandafter\def\csname xnumber@pp:prop:explicit\endcsname{7.1}
\expandafter\def\csname xnumber@pp:prop:fermatbcminimum\endcsname{3.1}
\expandafter\def\csname xnumber@pp:prop:fermatminpolarization\endcsname{3.2}
\expandafter\def\csname xnumber@pp:prop:fermatrestricted\endcsname{2.2}
\expandafter\def\csname xnumber@pp:prop:homideal\endcsname{7.8}
\expandafter\def\csname xnumber@pp:prop:idealclassmonoid\endcsname{7.10}
\expandafter\def\csname xnumber@pp:prop:markeddistinct\endcsname{7.3}
\expandafter\def\csname xnumber@pp:prop:markedpolarizations\endcsname{7.11}
\expandafter\def\csname xnumber@pp:prop:mindegree\endcsname{7.7}
\expandafter\def\csname xnumber@pp:prop:minimumpolarizedgalois\endcsname{5.2}
\expandafter\def\csname xnumber@pp:prop:picardorder\endcsname{7.9}
\expandafter\def\csname xnumber@pp:prop:polarizationsquare\endcsname{8.5}
\expandafter\def\csname xnumber@pp:prop:sixprincipal\endcsname{8.3}
\expandafter\def\csname xnumber@pp:prop:stabilizer\endcsname{7.2}
\expandafter\def\csname xnumber@pp:prop:trigonaltheta\endcsname{12.3}
\expandafter\def\csname xnumber@pp:prop:twoprimefermat\endcsname{11.1}
\expandafter\def\csname xnumber@pp:prop:unitcurve\endcsname{8.2}
\expandafter\def\csname xnumber@pp:prop:unitfamily\endcsname{8.1}
\expandafter\def\csname xnumber@pp:prop:unmarkeda\endcsname{8.4}
\expandafter\def\csname xnumber@pp:prop:unmarkeddistinct\endcsname{7.4}
\expandafter\def\csname xnumber@pp:prop:unpolarizeddistinct\endcsname{7.6}
\expandafter\def\csname xnumber@pp:rem:obstruction\endcsname{8.8}
\expandafter\def\csname xnumber@pp:rem:unitimage\endcsname{8.6}
\expandafter\def\csname xnumber@pp:sec:fermatintegral\endcsname{2}
\expandafter\def\csname xnumber@pp:sec:thetacertificate\endcsname{12.1}
\expandafter\def\csname xnumber@pp:thm:completequotientmoduli\endcsname{6.4}
\expandafter\def\csname xnumber@pp:thm:completequotientorbits\endcsname{6.2}
\expandafter\def\csname xnumber@pp:thm:definitiondegrees\endcsname{9.1}
\expandafter\def\csname xnumber@pp:thm:fermatallquotients\endcsname{4.1}
\expandafter\def\csname xnumber@pp:thm:fermatunderlying\endcsname{4.3}
\expandafter\def\csname xnumber@pp:thm:introgeometry\endcsname{1.3}
\expandafter\def\csname xnumber@pp:thm:jacobiancomplete\endcsname{12.4}
\expandafter\def\csname xnumber@pp:thm:obstructedisogenies\endcsname{10.2}
\expandafter\def\csname xnumber@pp:thm:obstructedminimum\endcsname{5.1}
\expandafter\def\csname xnumber@pp:thm:quotientdefinitionfield\endcsname{8.1}
\expandafter\def\csname xnumber@pp:thm:quotientgalois\endcsname{4.2}
\expandafter\def\csname xnumber@pp:thm:quotientrayfield\endcsname{6.1}
\expandafter\def\csname xnumber@pp:thm:quotientrealmodels\endcsname{7.1}
\expandafter\def\csname xnumber@pp:thm:realdefinitionfield\endcsname{8.2}
\expandafter\def\csname xnumber@pp:thm:tenprincipal\endcsname{12.1}
\expandafter\def\csname xnumber@prop:2816\endcsname{12.10}
\expandafter\def\csname xnumber@prop:Amaximalpair\endcsname{11.7}
\expandafter\def\csname xnumber@prop:Appisogenous\endcsname{10.3}
\expandafter\def\csname xnumber@prop:BCdistinct\endcsname{5.16}
\expandafter\def\csname xnumber@prop:Cppisogenous\endcsname{10.1}
\expandafter\def\csname xnumber@prop:Cvertex\endcsname{5.15}
\expandafter\def\csname xnumber@prop:DBsame\endcsname{5.17}
\expandafter\def\csname xnumber@prop:E1type\endcsname{4.1}
\expandafter\def\csname xnumber@prop:Kcompatible\endcsname{4.7}
\expandafter\def\csname xnumber@prop:actualcmtype\endcsname{5.7}
\expandafter\def\csname xnumber@prop:alltraceorders\endcsname{8.3}
\expandafter\def\csname xnumber@prop:broadersevenexclusions\endcsname{9.14}
\expandafter\def\csname xnumber@prop:canonical-fixed\endcsname{6.1}
\expandafter\def\csname xnumber@prop:charpoly\endcsname{12.7}
\expandafter\def\csname xnumber@prop:cmrepair\endcsname{12.8}
\expandafter\def\csname xnumber@prop:commonprincipal\endcsname{6.1}
\expandafter\def\csname xnumber@prop:d3ambientpoints\endcsname{3.2}
\expandafter\def\csname xnumber@prop:d7batch10\endcsname{9.5}
\expandafter\def\csname xnumber@prop:d7finalpairs\endcsname{9.6}
\expandafter\def\csname xnumber@prop:d7mask14\endcsname{9.1}
\expandafter\def\csname xnumber@prop:d7mask15_31\endcsname{9.2}
\expandafter\def\csname xnumber@prop:d7mask1_17\endcsname{9.4}
\expandafter\def\csname xnumber@prop:d7mask6_22\endcsname{9.3}
\expandafter\def\csname xnumber@prop:degree-obstruction\endcsname{12.3}
\expandafter\def\csname xnumber@prop:discminusone\endcsname{5.8}
\expandafter\def\csname xnumber@prop:effectiveorbit\endcsname{7.2}
\expandafter\def\csname xnumber@prop:fixedmatrixlattice\endcsname{5.4}
\expandafter\def\csname xnumber@prop:geometricB\endcsname{5.2}
\expandafter\def\csname xnumber@prop:knownCM\endcsname{5.18}
\expandafter\def\csname xnumber@prop:mainarith\endcsname{5.1}
\expandafter\def\csname xnumber@prop:noprincipal\endcsname{4.2}
\expandafter\def\csname xnumber@prop:orientationt\endcsname{7.1}
\expandafter\def\csname xnumber@prop:pi1\endcsname{3.1}
\expandafter\def\csname xnumber@prop:ppisogenous\endcsname{4.3}
\expandafter\def\csname xnumber@prop:principalstabilizer\endcsname{6.3}
\expandafter\def\csname xnumber@prop:residualexclusions\endcsname{11.9}
\expandafter\def\csname xnumber@prop:sixHG\endcsname{12.11}
\expandafter\def\csname xnumber@prop:sixvertexorbits\endcsname{11.6}
\expandafter\def\csname xnumber@prop:survivingorder\endcsname{6.5}
\expandafter\def\csname xnumber@q:jacobians\endcsname{12}
\expandafter\def\csname xnumber@rem:ABprincipallevels\endcsname{6.2}
\expandafter\def\csname xnumber@rem:WKnormalization\endcsname{1.4}
\expandafter\def\csname xnumber@rem:fig5signs\endcsname{5.6}
\expandafter\def\csname xnumber@rem:notgeometry\endcsname{5.2}
\expandafter\def\csname xnumber@rem:ppcover\endcsname{4.5}
\expandafter\def\csname xnumber@rem:pphec\endcsname{4.6}
\expandafter\def\csname xnumber@sec:actualcmtype\endcsname{5.4}
\expandafter\def\csname xnumber@sec:conclusion\endcsname{9}
\expandafter\def\csname xnumber@sec:fullhodge\endcsname{5.6}
\expandafter\def\csname xnumber@sec:hypergeom\endcsname{12.6}
\expandafter\def\csname xnumber@sec:intro\endcsname{1}
\expandafter\def\csname xnumber@sec:latticetransport\endcsname{5.3}
\expandafter\def\csname xnumber@sec:maincertificate\endcsname{5.1}
\expandafter\def\csname xnumber@sec:markman\endcsname{5.5}
\expandafter\def\csname xnumber@sec:polarization\endcsname{4}
\expandafter\def\csname xnumber@sec:setup\endcsname{2}
\expandafter\def\csname xnumber@sec:vertexhit\endcsname{5.2}
\expandafter\def\csname xnumber@six:sec:main\endcsname{3}
\expandafter\def\csname xnumber@six:tab:types\endcsname{1}
\expandafter\def\csname xnumber@six:thm:products\endcsname{3.1}
\expandafter\def\csname xnumber@subsec:arithmeticcensus\endcsname{3.1}
\expandafter\def\csname xnumber@tab:audit\endcsname{1}
\expandafter\def\csname xnumber@thm:42\endcsname{4.1}
\expandafter\def\csname xnumber@thm:A\endcsname{1.2}
\expandafter\def\csname xnumber@thm:Ahodge\endcsname{5.14}
\expandafter\def\csname xnumber@thm:Avertex\endcsname{5.13}
\expandafter\def\csname xnumber@thm:B\endcsname{12.1}
\expandafter\def\csname xnumber@thm:Bvertex\endcsname{5.3}
\expandafter\def\csname xnumber@thm:GLcollapse\endcsname{11.1}
\expandafter\def\csname xnumber@thm:d3classes\endcsname{8.6}
\expandafter\def\csname xnumber@thm:d7ambient\endcsname{8.7}
\expandafter\def\csname xnumber@thm:d7nonmax\endcsname{8.5}
\expandafter\def\csname xnumber@thm:d7r41complete\endcsname{9.11}
\expandafter\def\csname xnumber@thm:d7r83principal\endcsname{9.8}
\expandafter\def\csname xnumber@thm:fullhodge\endcsname{5.11}
\expandafter\def\csname xnumber@thm:fullpacketarithmetic\endcsname{3.3}
\expandafter\def\csname xnumber@thm:fullpacketgeometry\endcsname{3.6}
\expandafter\def\csname xnumber@thm:maximalallorientations\endcsname{9.10}
\expandafter\def\csname xnumber@thm:maxorder\endcsname{8.1}
\expandafter\def\csname xnumber@thm:packetHodge\endcsname{11.8}
\expandafter\def\csname xnumber@thm:packetintro\endcsname{1.1}
\expandafter\def\csname xnumber@thm:principalfamily\endcsname{6.4}
\expandafter\def\csname xnumber@thm:signature\endcsname{5.1}
\expandafter\def\csname xnumber@thm:unitnorm\endcsname{11.3}

\expandafter\def\csname doccode@packet\endcsname{I}
\expandafter\def\csname doccode@quotients\endcsname{II}
\expandafter\def\csname doccode@hodge\endcsname{H}
\expandafter\def\csname doccode@supplement\endcsname{S}

\title{Polarizations and Descent\\in a Finite Family of CM Sixfolds}
\author{Rollins Waltchack}
\date{3 November 2026}
\begin{document}
\maketitle
\begin{abstract}
Starting from the primitive six-dimensional factor of the Jacobian
of $y^2=x^{21}-1$, we classify a finite family of integral quotients and
minimum polarizations. The source and its order-$27$ quotients give $189$
underlying complex isomorphism classes and $2969$ minimum polarized
classes. Of the $188$ quotient classes, $80$ are principalizable and $108$
are not; we determine their minimum forms and distinguish target counts
from kernel counts. We compute the Galois action and fields of moduli,
then prove a descent obstruction and construct a common field of
definition of degree $24$. Minimum absolute definition degrees are sharp
for $168$ quotient classes, with $20$ cases left as lower bounds.
A restricted ten-class Jacobian classification is stated under its
explicit enclosure and execution assumptions; $156$ smaller-order
principal Jacobian questions remain open.
\end{abstract}
\paragraph{Companion references.}
Result numbers prefixed I, II, H and S refer respectively to the packet
paper, quotient paper, supporting Hodge note and computational supplement.
Each document has its own numbering and bibliography.
{\small\tableofcontents}
\clearpage
% !TeX root = ../../quotients-and-descent.tex
\section{Introduction}
This paper asks an integral question within one complex CM isogeny class:
which quotients of a fixed sixfold admit principal polarizations, and over
which fields can the quotients and their polarizations be defined? A
rational Hodge-algebraicity theorem does not answer either question.
The starting variety is defined from a curve, independently of any
Hilbert-modular packet.

Let $H$ be the smooth projective model of $y^2=x^{21}-1$. Its primitive
$\Phi_{21}$ factor $A\subset J(H)$ is six-dimensional and has full
endomorphism ring $\mathcal O_M$, where $M=\mathbb Q(\zeta_{21})$.
We keep this actual embedded factor distinct from isogenous principal
models. Its restricted Jacobian polarization has degree $35721$, whereas
its minimum polarization has degree $729$ and type $(1,1,1,3,3,3)$.

\subsection{The finite family and the equivalence relations}
Fix the minimum polarization $\lambda_{\min}$ and write its kernel as
$A[\pi]\simeq\mathbb F_3^6$, with the notation made explicit below.
The family consists of $A$ and all quotients $A/H_0$ for subgroups
$H_0\subset A[\pi]$ of order $27$. All isomorphisms preserve the origin.
An underlying class forgets the polarization; a polarized class retains
it but forgets CM and RM markings. A kernel is a subgroup of the fixed
source, not a target isomorphism class. These distinctions account for the
different counts in the following summary.

\begin{center}
\begin{tabular}{lr}
\toprule
Object being counted & Number\\\midrule
Order-$27$ subgroups of the fixed source kernel & $33880$\\
Underlying quotient classes (excluding $A$) & $188$\\
Underlying classes including $A$ & $189$\\
Principalizable underlying quotients & $80$\\
Principal polarized quotient classes & $160$\\
Nonprincipalizable underlying quotients & $108$\\
Minimum polarized classes including $A$ & $2969$\\\bottomrule
\end{tabular}
\end{center}
The last number is $1+160+108\cdot26$. The source contributes its unique
minimum polarization, and each obstructed quotient contributes $26$.
Fixing the pullback form gives $1120$ principal-target kernels; allowing
all minimum pullback forms gives $14224$. Neither number counts all
$33880$ order-$27$ subgroups as principal-target kernels.

\subsection{Main results and proof route}
Theorem~\ref{pp:thm:fermatallquotients} classifies the principal targets;
Theorem~\ref{pp:thm:fermatunderlying} forgets the polarizations and identifies
the obstruction; Theorem~\ref{pp:thm:obstructedminimum} determines the
minimum forms on the remaining $108$ varieties. Their proofs pass from
integral lattices to finite residue subspaces and distinguish the relevant
unit actions before taking orbits.

The Galois action then determines fields of moduli, but a fixed isomorphism
class need not descend. Put
\[
 \mathcal E=M\bigl(\sqrt{(3+\sqrt{21})/2}\bigr).
\]
Theorem~\ref{pp:thm:quotientdefinitionfield} proves that for number fields $M\subseteq F\subset\overline{\mathbb Q}$, any one of the $188$ quotient classes has an $F$-model
exactly when $F$ contains $\mathcal E$. A single twist makes all subgroup
kernels and quotient maps available over $\mathcal E$. Thus the common
field-of-moduli computation and the actual defining-field construction
are distinct stages of the argument.
Theorem~\ref{pp:thm:definitiondegrees} gives sharp absolute degrees for
$168$ quotient classes and lower bounds for the other $20$.

The final geometric section treats Jacobians. It distinguishes the
restricted ten-class principal census with $R$-containment from the $156$
principal classes with smaller endomorphism order $S$. The former has
exactly one canonical Jacobian under the stated enclosure and execution
contract. The latter $156$ statuses remain open.

\subsection{Arithmetic inputs and relation to other work}
Write $z=\zeta_{42}=-\zeta_{21}^{11}$, $t=z+z^{-1}$, $L=M^+$ and
\[
 f(T)=T^6+T^5-6T^4-6T^3+8T^2+8T+1,\qquad
 \mathfrak P=(1-z^{14})\mathcal O_M,
 \quad R=\mathcal O_L+\mathfrak P,\quad S=\mathbb Z+\mathfrak P.
\]
Here $\mathfrak P$ is the prime over $3$; its residue field is
$\mathbb F_{729}$ and $R/\mathfrak P=\mathbb F_{27}$. The letter $\pi$
below is $1-z^{14}$, a generator of $\mathfrak P$ in $M$. It is \emph{not} the norm-$43$ generator of
Paper I. We use $K=\mathbb Q(\sqrt{-3})$ for the Weil action.
Write $E_L=\mathcal O_L^\times$ and let $E_L^+$ denote its subgroup
of totally positive units.

The imported arithmetic inputs are the ideal classification and marked
polarization classification of Paper I, Propositions~\externalref{packet}{pp:prop:picardorder},
\externalref{packet}{pp:prop:idealclassmonoid} and~\externalref{packet}{pp:prop:markedpolarizations},
and the unit-norm theorem in the supplement. They are established without
using this quotient family. The notations $B_0,B_1,C_0,C_1$ mean the two
principal-polarization classes on each of the two $R$-ideal models in the
B/C CM-type orbit; Paper I identifies them with its principal curve fibers.
The quotient construction here needs their arithmetic classification, not
that curve placement. The type representatives are
$\Phi_B=\{1,5,23,25,29,31\}$ and
$\Phi_C=\{1,5,11,13,17,19\}$; $\Phi_0=\Phi_C$.

Goodson and Saller discuss the rational CM-type/Fermat comparison
\cite{Goodson2024,Saller2026}. The general CM Hodge assertion in
\cite{OpenAICM2026}, if valid, concerns algebraicity of rational classes;
it does not itself provide the integral census or descent data here.
The existence of individual models must also be compared with classical
CM descent, including Milne's correction and addendum
\cite{MilneModuli1972,MilneModuli1974}. Our explicit simultaneous twist,
field obstruction and finite tables are distinguished from that existence
question. These comparisons are not a certification of historical priority.

% !TeX root = ../../quotients-and-descent.tex
\section{The source factor and its polarization}
\label{pp:sec:fermatintegral}
Let $M=\mathbf Q(\zeta_{21})=\mathbf Q(\zeta_{42})$, with
$\zeta_{42}=-\zeta_{21}^{11}$. The following construction separates a
rational CM-type comparison from the principal polarizations classified in Section~\ref{q:jacobians}.
The order-and-type census used here is proved in Theorem~\ref{pp:thm:tenprincipal}; its proof uses only the imported order arithmetic, not the quotient or descent conclusions.

\begin{proposition}
Let $H$ be the smooth projective curve with affine equation
$y^2=x^{21}-1$. Put $\sigma(x,y)=(\zeta_{21}x,y)$ and let
$\sigma_J([p-\infty])=[\sigma(p)-\infty]$ on $J(H)$.
The connected image $A$ of a denominator-cleared primitive
$\Phi_{21}$ idempotent is a simple abelian sixfold with
$\operatorname{End}(A)=\mathcal O_M$. Its rational CM type belongs to the
B/C orbit, but $A$ admits no principal polarization. In particular, the
polarization restricted from $J(H)$ is nonprincipal.
\end{proposition}

\begin{proof}
The double cover of the $x$-line has twenty-two simple branch points,
including infinity, so $H$ has genus ten. A basis of holomorphic forms is
$x^{i-1}dx/y$, $1\leq i\leq10$. At infinity the orders of $x,y,dx$ are
$-2,-21,-3$, respectively, giving order $20-2i\geq0$ for these forms;
at a finite branch point $dx/y$ is regular in the parameter $y$.
The eigenvalue of $\sigma^*$ on the $i$th form is $\zeta_{21}^i$.
The primitive characters are exactly
\[
T=\{1,2,4,5,8,10\}.
\]
Their complex conjugates complete the twelve primitive characters in
rational cohomology. Abel--Jacobi equivariance identifies these with the
cotangent eigenvalues of the stated $\sigma_J$; the derivative on the tangent space and its dual on cotangent
forms have the same eigenvalues. This convention is essential:
functorial pullback on $\operatorname{Pic}^0(H)$ would give $\sigma_J^{-1}$.

The Chinese remainder decomposition of $\mathbf Q[t]/(t^{21}-1)$ gives
a rational idempotent $e$ for $\Phi_{21}$, acting by $t=\sigma_J$.
Choose an integer $N>0$ such that $Ne$ is an integral endomorphism.
The image of this abelian-variety homomorphism is a connected abelian
subvariety; its tangent space is the image of the derivative. Thus
$A=\operatorname{im}(Ne)$ has dimension six. Its image does not depend on
$N$, since multiplication by a nonzero integer is surjective on an
abelian variety. The tangent CM type for $\zeta_{21}\mapsto\sigma_J|_A$
is $T$. Under $\zeta_{42}=-\zeta_{21}^{11}$ its exponents become
$\{1,5,23,25,29,31\}$, the actual B type.

Both $\sigma_J$ and its inverse preserve $A$, since they commute with $e$.
Their restriction is an integral automorphism with minimal polynomial
$\Phi_{21}$, so
$\mathcal O_M=\mathbf Z[\zeta_{21}]\subseteq\operatorname{End}(A)$.
The primitive-type endomorphism calculation in Theorem~\ref{pp:thm:tenprincipal} gives simplicity and
$\operatorname{End}^0(A)=M$ without requiring a principal polarization.
Maximality of $\mathcal O_M$ then gives equality integrally.
If $A$ admitted a principal polarization, it would occur in Theorem~\ref{pp:thm:tenprincipal}'s complete
census: its endomorphisms contain the specified order $R$, and the Rosati
involution on the CM field is complex conjugation. But the B/C type orbit
has no principal object with full endomorphism order $\mathcal O_M$.
This contradiction proves the assertion. Restriction of the ample
Jacobian polarization is a polarization on $A$, hence is nonprincipal.
\end{proof}

There is also a direct Fermat domination. For $c^{21}=4$, the formula
\[
 (u,v)\longmapsto (X,Y)=(cuv,\,\mathrm i(u^{21}-v^{21}))
\]
on $u^{21}+v^{21}=1$ satisfies $Y^2=X^{21}-1$.
The nonconstant map extends to the smooth projective models. The identity
$f_*f^*=[\deg f]$ on Jacobians exhibits $J(H)$, and therefore $A$, as
isogeny factors of the Fermat Jacobian.
This does not identify $A$ with any of the four principally polarized
B/C objects: their full endomorphism order is $R$, whereas that of $A$
is $\mathcal O_M$. Proposition~\ref{pp:prop:fermatrestricted} below
computes the actual restricted polarization. Proposition~\ref{pp:prop:fermatbcminimum}
determines the minimum unpolarized isogenies to B/C. A polarized identification
with a particular prior-art factor remains a separate question. Their genus-six Jacobian status is instead addressed by
Theorem~\ref{pp:thm:jacobiancomplete}, under its stated computational contract.

\begin{proposition}[The induced polarization on the Fermat factor]
\label{pp:prop:fermatrestricted}
For the actual embedded primitive factor $A\subset J(H)$ above,
the restriction of the canonical principal polarization has type
$(1,1,1,3,3,21)$ and degree $35721$. Put
$H_d:y^2=x^d-1$ for $d=3,7$, taking smooth projective models,
and let $f_d:H\to H_d$ be $(x,y)\mapsto(x^{21/d},y)$.
The addition map
\[
 A\times J(H_3)\times J(H_7)\longrightarrow J(H),\qquad
 (a,u,v)\longmapsto a+f_3^*u+f_7^*v
\]
is an isogeny of degree $35721$.
\end{proposition}
\begin{proof}
The lower curves have genera one and three. Each $f_d$ has prime
degree $q=21/d$ and is ramified: at $(0,\pm\mathrm i)$ the local
parameter $x$ maps to $x^q$. The norm identity gives
$f_{d,*}f_d^*=[q]$. If $\ker f_d^*$ contained a nonzero line bundle
$L$, it would have exact order $q$. Its connected cyclic \'etale
cover becomes trivial after pullback to $H$, so $f_d$ would factor
through that cover. The remaining map would have degree one,
contradicting ramification. Thus $f_d^*$ is an embedding.
The canonical polarization adjunction, equivalently the degree
formula for the integral cup pairing, gives induced form
$q\lambda_{H_d}$ on its image $Z_d$.

Let $T=\sigma_J$ in the covariant convention above, with compatible
roots $\zeta_d=\zeta_{21}^{21/d}$. Equivariance
implies that $\Phi_d(T)$ annihilates $Z_d$: the sum of the
order-$d$ rotations factors through the Jacobian of the quotient
$y$-line, which is zero. The integral identity
\[
 \Phi_7(t)-(t+t^4)\Phi_3(t)=1
\]
therefore proves $Z_3\cap Z_7=0$. Since $T$ preserves the canonical
alternating form, eigenvectors pair to zero unless their eigenvalues
are inverse. Primitive third and seventh characters are orthogonal.
Thus $Z=Z_3+Z_7$ is the embedded product, with restricted form
$7\lambda_{H_3}\oplus3\lambda_{H_7}$. Its polarization kernel is
\[
 (\mathbb Z/7)^2\oplus(\mathbb Z/3)^6
 \simeq(\mathbb Z/3)^4\oplus(\mathbb Z/21)^2,
\]
so its type is $(1,3,3,21)$.

The remaining rational homology characters are exactly the primitive
twenty-first characters. They identify the actual $A$ with the
connected orthogonal complement of $Z$. To keep the integral step
explicit, let $\Lambda=H_1(J(H),\mathbb Z)$ and let $M,N$ be its
saturated sublattices in the rational spaces of $Z,A$. The ambient
form is unimodular. Projection of $\Lambda$ onto $M_{\mathbb Q}$
maps onto the dual lattice $M^\#$: every integral functional on the
primitive sublattice $M$ extends to $\Lambda$, and unimodularity
realizes it by pairing with a lattice element. The projection kernel
is $N$, whence
\[
 \Lambda/(M\oplus N)\simeq M^\#/M\simeq N^\#/N.
\]
The two restrictions have the same nonunit elementary divisors.
The six-dimensional complement therefore adds two unit entries to
the type of $Z$, giving $(1,1,1,3,3,21)$. The addition kernel is the
displayed quotient, of order $7^2 3^6=35721$, which is also the degree
of the restricted polarization on $A$.
\end{proof}

The degrees here are degrees of the specified polarization and
addition map; the product of the six polarization elementary divisors
is $189$. The minimum unpolarized B/C isogeny degree is instead $27$
by Proposition~\ref{pp:prop:fermatbcminimum}. This structural argument uses neither the
theta computation nor its analytic-library contract.

\begin{corollary}[An integral primitive projector]
\label{pp:cor:fermatprojector}
With $T=\sigma_J$, the endomorphism
\[
 \mathcal N=21-3\sum_{j=0}^{6}T^{3j}-7(1+T^7+T^{14})
\]
has image $A$, is symmetric for the Rosati involution of $J(H)$,
and satisfies $\mathcal N^2=21\mathcal N$.
For the primitive projector $e=\mathcal N/21$, the least positive
integer $n$ such that $ne\in\operatorname{End}(J(H))$ is $21$.
\end{corollary}
\begin{proof}
The deck-group identities for the actual covers give
\[
 f_3^*f_{3,*}=\sum_{j=0}^{6}T^{3j},\qquad
 f_7^*f_{7,*}=1+T^7+T^{14}.
\]
The first sum acts by $7$ on the order-three characters and by zero
on the others; the second acts by $3$ on the order-seven characters
and by zero on the others. Thus $\mathcal N/21$ is precisely the
orthogonal projector onto the primitive part. Multiplication by $21$
is surjective on $A$, proving the image assertion. Rosati sends $T$
to $T^{-1}$ and preserves both subgroup sums. This proves symmetry;
the projector identity gives $\mathcal N^2=21\mathcal N$.
Equivalently, the polynomial residual is
$42(1+T+\cdots+T^{20})$, which vanishes on $J(H)$ because the full
rotation quotient is the projective $y$-line.

Write $L_A=H_1(A,\mathbb Z)$ and $\Lambda=H_1(J(H),\mathbb Z)$.
The saturated-lattice projection argument in
Proposition~\ref{pp:prop:fermatrestricted} gives $e\Lambda=L_A^\#$.
Consequently $ne$ preserves $\Lambda$ exactly when
$nL_A^\#\subseteq L_A$. This is equivalent to $n$ annihilating
$L_A^\#/L_A$, whose exponent is $21$. Lattice preservation gives a
holomorphic endomorphism of the complex torus, hence an algebraic
endomorphism of this abelian variety. Therefore $21$ is necessary
and sufficient.
\end{proof}

This is minimality of the denominator of this specified projector;
it is not a minimal isogeny statement for the principal B/C objects.


\section{Minimum isogenies and the minimum source polarization}
\begin{proposition}[Minimum isogenies from the Fermat factor]
\label{pp:prop:fermatbcminimum}
Let $A$ be the actual primitive factor of $J(H)$ above, and let
$B,C$ denote the two underlying sixfolds of the four principal objects
$B_0,B_1,C_0,C_1$. Among origin-preserving complex unpolarized
isogenies, the minimum degree in each of the directions
$A\to B$, $B\to A$, $A\to C$, $C\to A$ is $27$.
Every minimum kernel is isomorphic to $(\mathbb Z/3)^3$.

For the fixed source $A$, the minimum kernels giving a quotient
isomorphic to $B$ or $C$ are exactly the $28$ one-dimensional
$\mathbb F_{27}$-subspaces of
$A[\pi]\simeq\mathcal O_M/\mathfrak P=\mathbb F_{729}$,
where $\pi=1-z^{14}$: fourteen give each underlying target.
For each fixed source $B$ or $C$, there is exactly one minimum
kernel giving a quotient isomorphic to $A$.
\end{proposition}
\begin{proof}
Put $O=\mathcal O_M$. Since $\operatorname{End}(A)=O$ and
$\operatorname{Pic}(O)=1$, its actual rank-one CM lattice can be
scaled to $O$. At a common CM type, the target lattice can be
scaled to a line lift $I$ with
\[
 \mathfrak P\subset I\subset O,\qquad IO=O,\qquad [O:I]=27.
\]
These are isomorphisms of complex tori. The needed field re-marking
preserves $O,R,\mathfrak P$ and the two classes of
$\operatorname{Pic}(R)\simeq\mathbb Z/2$.
Thus the actual $B,C$ classes remain $[\mathfrak a],[R]$, respectively, by
Propositions~\externalref{packet}{pp:prop:picardorder},
\externalref{packet}{pp:prop:idealclassmonoid} and \externalref{packet}{pp:prop:actualpolarizations}.

Reduction modulo $\mathfrak P$ proves
$(I:O)=\mathfrak P$: a nonzero residue scalar cannot send the
whole field $\mathbb F_{729}$ into a proper
$\mathbb F_{27}$-line. The equality $IO=O$ gives $(O:I)=O$.
Multiplication by $\pi$ and inclusion supply isogenies in the two
directions. Every other unmarked isogeny is scalar in these common
coordinates, since composing with a rational inverse of either one
lands in the full endomorphism algebra $M$. Therefore their degrees
are exhausted by
\[
 [I:cO]=\frac{|N_{M/\mathbb Q}(c)|}{27}
 \quad(c\in\mathfrak P\setminus\{0\}),\qquad
 [O:cI]=27|N_{M/\mathbb Q}(c)|
 \quad(c\in O\setminus\{0\}).
\]
The first minimum is $729/27=27$, attained exactly for
$c=\pi u$, $u\in O^\times$; the second is $27$, attained
exactly for $c\in O^\times$.

The forward kernel is $\pi^{-1}u^{-1}I/O$. Multiplication by
$\pi$ identifies it with the line $u_{\mathrm{res}}^{-1}(I/\mathfrak P)$
in $O/\mathfrak P$, where $u_{\mathrm{res}}$ is the residue of $u$.
The backward kernel is always
$O/I$. Both groups are one-dimensional over $\mathbb F_{27}$.
The $28$ projective lines form a torsor for
$\mathbb F_{729}^{\times}/\mathbb F_{27}^{\times}$; the full
global-unit image has order $14$ by
Proposition~\externalref{packet}{pp:prop:picardorder}. Its two orbits give precisely
the two ideal classes. Conversely every line lift $I'$ yields a
kernel $\pi^{-1}I'/O$, whose quotient is scalar-isomorphic to the
torus of lattice $I'$. This proves exhaustion and the split
$14+14$. The backward formula proves uniqueness for each fixed
source. Changing the scalar normalization permutes the same lines.
\end{proof}

These are counts of subgroups and minimum degrees of unpolarized
maps. They impose no compatibility with the specified polarizations
and do not specify a field of definition or a correspondence of curves.


\begin{proposition}[Minimum polarization and principal quotients]
\label{pp:prop:fermatminpolarization}
The minimum degree of a polarization on the actual primitive factor
$A\subset J(H)$ is $729$. There is exactly one such polarization
up to origin-preserving automorphism of $A$; its type is
$(1,1,1,3,3,3)$ and its kernel is $A[\pi]$.

Fix a representative $\lambda_{\min}$. Each of the $28$
$\mathbb F_{27}$-lines in $A[\pi]$ is maximal isotropic and gives
a principally polarized quotient with pullback $\lambda_{\min}$.
These quotients have exactly four polarized isomorphism classes,
each realized by seven distinct kernels: they are $B_0,B_1,C_0,C_1$.
\end{proposition}
\begin{proof}
Normalize the actual CM lattice of $A$ to $O=\mathcal O_M$ as in
Proposition~\ref{pp:prop:fermatbcminimum}. The rotation preserves the
restricted Jacobian polarization and generates $M$, so its Rosati
involution is complex conjugation. Any other polarization differs
from it by an element of the commutative algebra
$\operatorname{End}^0(A)=M$, and therefore induces the same Rosati
involution. Thus the trace-form classification applies to every
polarization, without an additional compatibility assumption.

Write $\kappa_O=1/(f'(t)(z-\bar z))$, so $O^\vee=\kappa_O O$.
Integral alternating forms have parameters $\xi=\kappa_O a$ with
$a\in\mathcal O_L$. At the actual B type, positivity forces the
sign vector of $a$ to equal that of $\operatorname{Im}\kappa_O$,
which has odd negative parity by
Proposition~\externalref{packet}{pp:prop:markedpolarizations}. Hence
$N_{L/\mathbb Q}(a)<0$, and the degree is
$N_{L/\mathbb Q}(a)^2$ because the dual lattice is $a^{-1}O$.

We claim that no real integer has negative norm of absolute value
less than $27$. Since $\mathcal O_L=\mathbb Z[t]$, the following
reductions of its defining polynomial determine the possible prime
ideal norms below $27$:
\[
\begin{array}{c|l}
p&f\bmod p\\ \hline
2&\text{irreducible of degree }6\\
3&(t^3-t^2+t+1)^2\\
5&(t^3-t-2)(t^3+t^2+2)\\
7&(t-1)^6.
\end{array}
\]
The displayed cubics are irreducible by absence of roots. Modulo
$2$, division by all fourteen monic polynomials of degrees at most
three verifies irreducibility. For $p=11,13,17,19,23$, direct
evaluation gives no root, so the prime ideal norms are at least
$p^2>27$. This exhausts all rational primes below $27$. The finite reductions,
root checks and norms here are reproduced exactly by
\texttt{certificates/fermat\_minimum\_polarization\_norm\_v154.py};
that certificate supplies arithmetic inputs to the geometric argument.
The only prime ideal of norm below $27$ is consequently the one
of norm $7$, generated by $t-1$, whose field norm is $f(1)=7$.
An ideal of norm below $27$ is either this ideal or the unit ideal.
All real units have norm $+1$ by Lemma~\externalref{supplement}{lem:narrowclass}, so
all their generators have positive norm, proving the claim.

Every polarization therefore has degree at least $729$. A minimum
degree-$27$ map of Proposition~\ref{pp:prop:fermatbcminimum} to
principally polarized B pulls back a polarization of degree $729$.
Every minimum parameter $a$ generates the unique real prime above
$3$, which has ideal norm $27$ and is generated by $\beta$ with
$N_{L/\mathbb Q}(\beta)=-27$. Relative to one
positive minimum parameter, all others differ by $E_L^+$.
Theorem~\externalref{supplement}{thm:unitnorm} gives
$N_{M/L}(O^\times)=E_L^+$, proving uniqueness up to automorphism.
Since $\beta O=\mathfrak P=\pi O$, the kernel is
$\pi^{-1}O/O\simeq(\mathbb Z/3)^6$, giving the stated type.

The discriminant pairing on this kernel is alternating over
$\mathbb F_3$ and self-adjoint for the real residue subfield
$k=\mathbb F_{27}$. Thus for $a\in k$,
$b(x,ax)=b(ax,x)=-b(x,ax)$, whence $b(x,ax)=0$.
Every $k$-line is therefore isotropic and, having order $27$,
maximal isotropic. Its inverse image in $O^\#$ is an integral
self-dual overlattice for the same positive Riemann form. This
gives the claimed principal quotient and pullback equality.

The polarized automorphism group of $A$ is $\mu_{42}$:
an $O$-unit preserving the form satisfies $u\bar u=1$ and is a
root of unity. Its image in
$\mathbb F_{729}^\times/\mathbb F_{27}^\times$ has order $7$,
so it acts freely with four seven-element orbits on the lines.
These orbits cannot merge under polarized quotient isomorphisms.
Indeed, the quotient lattices are $\pi^{-1}I$ and all have
$O$-saturation $\pi^{-1}O$. Every unmarked isomorphism between
them is scalar by $\operatorname{End}^0(A)=M$. Equality of the
saturations forces this scalar to be an $O$-unit, and preserving
their common rational Riemann form forces it into $\mu_{42}$.
Conversely every such root induces the corresponding isomorphism.
Each fourteen-line B/C orbit splits into two seven-line orbits.
Proposition~\externalref{packet}{pp:prop:actualpolarizations} identifies these as
the two principal classes on each underlying variety.
\end{proof}

The minimum polarization here differs from the degree-$35721$
restriction of Proposition~\ref{pp:prop:fermatrestricted}.
The quotient statement in Proposition~\ref{pp:prop:fermatminpolarization}
concerns the specified $\mathbb F_{27}$-lines; the
following theorem treats all maximal isotropic subgroups. Neither
statement specifies a field of definition or an explicit curve correspondence.


\section{Principal and underlying quotient classifications}
\begin{theorem}[All principal targets at minimum degree]
\label{pp:thm:fermatallquotients}
The minimum degree of an isogeny from the actual primitive factor $A$
to any principally polarized abelian variety is $27$. Exactly $160$
origin-preserving polarized isomorphism classes of principal targets
admit an isogeny of this degree from $A$.

Fix $\lambda_{\min}$ from Proposition~\ref{pp:prop:fermatminpolarization}.
The maps with pullback $\lambda_{\min}$ have exactly $1120$ distinct
kernels, the maximal isotropic subspaces of $A[\pi]\simeq\mathbb F_3^6$.
They give the $160$ polarized target classes with seven kernels each.
Four are $B_0,B_1,C_0,C_1$, with endomorphism ring $R$. The remaining
$156$ have full endomorphism ring
\[
 S=\mathbb Z+\mathfrak P\subset\mathcal O_M,\qquad
 [\mathcal O_M:S]=243.
\]
Each target has polarized automorphism group $\mu_6$.
\end{theorem}
\begin{proof}
A degree-$d$ map to a principal target pulls back a polarization of
degree $d^2$, so Proposition~\ref{pp:prop:fermatminpolarization}
gives $d\ge27$, attained by the B/C maps. For any degree-$27$ map,
uniqueness of the minimum form gives a source automorphism transporting
its pullback to $\lambda_{\min}$. This preserves the polarized target,
so it suffices to fix that pullback when counting target classes.

Use the actual lattice $O=\mathcal O_M$ and its minimum Riemann form
$E$. Its dual is $O^\#=\pi^{-1}O$. A subgroup
$H\subset D=O^\#/O$ gives an integral quotient lattice $L_H$ precisely
when it is isotropic, and $L_H^\#/O=H^\perp$. Thus the same positive
form is principal on the quotient precisely when $H=H^\perp$.
Conversely every map with fixed pullback gives such a self-dual
overlattice, so these are exactly its kernels.

If $L_n(q)$ counts Lagrangians in a symplectic space of dimension
$2n$ over $\mathbb F_q$, counting pairs $(v,H)$ with $0\ne v\in H$
gives
\[
 L_n(q)(q^n-1)=(q^{2n}-1)L_{n-1}(q),\qquad L_0(q)=1.
\]
Indeed those containing $v$ correspond to Lagrangians in
$v^\perp/\langle v\rangle$. Consequently
$L_3(3)=(3+1)(9+1)(27+1)=1120$.

Under multiplication by $\pi$, view $H$ as a subspace of
$K_{\rm res}=O/\mathfrak P=\mathbb F_{729}$, and write
$L_H=\pi^{-1}I_H$, where $I_H$ is its inverse image in $O$.
Every nonzero $H$ gives $I_HO=O$, since its nonzero residues
generate the field $K_{\rm res}$ as an ideal. All quotient lattices
therefore have the same $O$-saturation $\pi^{-1}O$.
Every unmarked quotient isomorphism is scalar in $M$, by conjugating
with the rational quotient maps and using $\operatorname{End}^0(A)=M$.
Equality of saturations forces the scalar into $O^\times$; preserving
the common form $E$ forces $u\bar u=1$, hence $u\in\mu_{42}$.
Conversely such roots give the corresponding quotient isomorphisms.
The target classes are exactly the $\mu_{42}$-orbits of Lagrangians.

The residue image of $\mu_{42}$ has order $14$, with scalar
subgroup $\{\pm1\}$, so the effective action has order $7$.
A nonidentity effective element is represented, up to sign, by a
primitive seventh root. It generates $\mathbb F_{729}$ over
$\mathbb F_3$, since $\operatorname{ord}_7(3)=6$.
It cannot preserve a three-dimensional subspace: such a subspace
would then be a nonzero proper vector subspace over the whole field.
The action is free, giving $1120/7=160$ classes, and each full
stabilizer is the six-element scalar kernel $\mu_6$.

Finally, a multiplier $cI_H\subseteq I_H$ lies in $O$ after
extending to $O$. Its residue must belong to
\[
 F_H=\{a\in K_{\rm res}:aH\subseteq H\}.
\]
This is a subfield: a nonzero member acts bijectively on finite $H$,
so its inverse also preserves $H$. Its degree over $\mathbb F_3$
divides $\dim_{\mathbb F_3}H=3$. Hence $F_H$ is either
$\mathbb F_3$ or $\mathbb F_{27}$, the latter exactly for the
$28$ one-dimensional $\mathbb F_{27}$-subspaces. Their four
polarized classes are the B/C objects already classified.
All other $1092$ Lagrangians have multiplier equal to the inverse
image of $\mathbb F_3$, namely $S=\mathbb Z+\mathfrak P$,
of index $729/3=243$. They give $1092/7=156$ classes.
Full rational endomorphisms remain $M$ by isogeny, so these
multiplier orders are the full integral endomorphism rings.
\end{proof}

The $1120$ kernel count fixes the pullback polarization; the $160$
target count allows arbitrary minimum pullbacks. The smaller order
$S$ contains $\mathbb Z[z^{14}]=\mathcal O_{\mathbb Q(\sqrt{-3})}$,
since $z^{14}\equiv1\pmod{\mathfrak P}$, but does not contain $R$.
Thus these $156$ objects lie outside the selected-order hypothesis
of Theorem~\ref{pp:thm:tenprincipal}. The census does not identify
new points on the normalized McMullen curve or classify the new
targets as Jacobians.



\begin{theorem}[Galois orbits of the minimum principal targets]
\label{pp:thm:quotientgalois}
The $160$ polarized classes of
Theorem~\ref{pp:thm:fermatallquotients} form eleven orbits under $\operatorname{Gal}(\overline{\mathbb Q}/\mathbb Q)$.
One consists of $B_0,B_1,C_0,C_1$. The other ten, all of full order $S$,
have sizes
\[
4,8,12,12,12,12,24,24,24,24.
\]
In particular, the property of being the canonically polarized Jacobian
of a smooth curve for the $156$
smaller-order classes can be decided on ten representatives.
\end{theorem}
\begin{proof}
First the literal primitive factor $A$ has a model over $\mathbb Q$.
The curve $y^2=x^{21}-1$ is defined over $\mathbb Q$, and Galois sends
its Jacobian rotation $T$ to $T^a$. The primitive cyclotomic projector
is invariant under this substitution, which permutes primitive
characters. After clearing denominators it is a Galois-invariant
homomorphism, hence descends to $\mathbb Q$, as does its connected
image. All endomorphisms of $A$ are defined over $M$, since $T$ generates
$\mathcal O_M$; their Galois action is the cyclotomic action.

Fix a minimum polarization $\lambda$. Every conjugate is another
minimum polarization on $A$, so
Proposition~\ref{pp:prop:fermatminpolarization} gives an isomorphism
$u_\sigma:(A,\sigma\lambda)\longrightarrow(A,\lambda)$, unique modulo
$\mu_{42}$. This does not assert that $\lambda$ descends to $\mathbb Q$.
The maps $u_\sigma\sigma$ act on
$D=\ker\lambda\simeq K_{\rm res}=\mathbb F_{729}$ modulo the residue group
$\mu_{14}$. They induce the actual Galois action on the polarized
quotient classes: quotients and their polarizations commute with base
change, and Theorem~\ref{pp:thm:fermatallquotients} identifies all
polarized quotient isomorphisms with source root actions.

The discriminant pairing has the form
\[
b(x,y)=\operatorname{Tr}_{K_{\rm res}/\mathbb F_3}(cxy^{27}),
\qquad c^{27}=-c\ne0.
\]
This follows from trace nondegeneracy, adjoint conjugation and
alternation. Its semilinear similitude group is
\[
G=\{x\mapsto ux^{3^j}:u^{28}=\varepsilon c^{3^j-1},\quad
               j\in\mathbb Z/6,\ \varepsilon\in\mathbb F_3^\times\}.
\]
The norm $K_{\rm res}^\times\to\mathbb F_{27}^\times$ is surjective with
kernel $C_{28}$, giving an exact sequence
\[
1\longrightarrow C_{28}\longrightarrow G
 \longrightarrow C_6\times C_2\longrightarrow1.
\]
Thus $\overline G=G/\mu_{14}$ has order $24$.
For $\sigma(\zeta_{42})=\zeta_{42}^a$, the normalized action has
field exponent $3^j\equiv a\pmod7$ and pairing multiplier
$\varepsilon\equiv a\pmod3$, by Galois functoriality of the
polarization pairing. Consequently its image in $\overline G$
surjects onto $C_6\times C_2$.

To obtain the remaining kernel of order two, use CM reciprocity on
the maximal-order source. The actual type
$\Phi=\{1,5,11,13,17,19\}$ has trivial stabilizer; its reflex field
is $M$ and its reflex type is the inverse type
\cite[Example~1.19, Proposition~1.26]{MilneCM2020}.
For an idele $s$ supported in the local units at $\mathfrak P$,
with residue $r\in K_{\rm res}^\times$, the reflex norm has residue
\[
r^{1+3+9+27+243+3}=r^{286}.
\]
Since $\gcd(286,728)=26$, these residues exhaust $C_{28}$.
Choose $\sigma$ fixing $M$ with
$\sigma|_{M^{\rm ab}}=\operatorname{art}_M(s)$, using Milne's
reciprocal Artin convention.
The fundamental theorem supplies an $M$-quasi-isogeny
$\alpha:A\to\sigma A$ with $\alpha N_\Phi(s)=\sigma$ on the finite
rational Tate module \cite[Theorem~9.10]{MilneCM2020}.
The reflex norm is an adelic unit, so $\alpha$ and its inverse are
integral on every Tate module and therefore are actual isomorphisms.
The polarization multiplier is
$\chi_{\rm cyc}(\sigma)/N_{M/\mathbb Q}(s)$
\cite[Remark~9.11(c)]{MilneCM2020}. By Lemma~9.5 of that source it is
a positive rational number; both numerator and denominator are
adelic units, so it equals one. Hence $\alpha$ is polarized.
Taking $u_\sigma=\alpha^{-1}$ realizes $C_{28}/\mu_{14}$ in the
normalized Galois image. That image is therefore all of $\overline G$.

The exact finite action on the adopted $1120$ isotropic subspaces,
grouped into $160$ root orbits, gives the eleven asserted orbits.
Certificate \texttt{s\_order\_galois\_v163} retains the actual
discriminant form and coordinate map; a separate implementation enumerates semilinear maps in those
coordinates and verifies the same partition. No numerical periods
enter this step. Finally, Galois conjugation preserves canonical
Jacobians of smooth curves, in both directions, proving the last
assertion.
\end{proof}

This orbit theorem does not decide any of the ten remaining Jacobian
questions. It permits a separately certified exclusion on one
representative to apply to its orbit; it does not identify theta values
or torsion labels between distinct normalized period matrices.


\begin{theorem}[Underlying quotients and the principal obstruction]
\label{pp:thm:fermatunderlying}
The $33880$ subgroups $H\subset A[\pi]$ of order $27$ give exactly
$188$ origin-preserving unpolarized isomorphism classes of quotients:
two with full endomorphism ring $R$, namely $B,C$, and $186$ with
full endomorphism ring $S=\mathbb Z+\mathfrak P$.
Exactly $80$ underlying varieties admit principal polarizations: $B,C$
and $78$ of the $S$ varieties. Each admits exactly two principal
polarization classes under its full automorphism group. The other
$108$ varieties admit no principal polarization.

Allowing arbitrary minimum pullback forms, the degree-$27$ isogenies
from fixed $A$ to principally polarizable targets have exactly $14224$
distinct kernels: $14$ for each $R$ target and $182$ for each of the
$78$ eligible $S$ targets. Fixing $\lambda_{\min}$ instead gives the
$1120$ kernels of Theorem~\ref{pp:thm:fermatallquotients}, with $14$
per underlying target, seven for each principal polarization class.
\end{theorem}
\begin{proof}
Put $k=\mathbb F_{27}\subset K_{\rm res}=\mathbb F_{729}$.
The real unit $t$ has residue with irreducible polynomial
$X^3-X^2+X+1$. Its norm to $\mathbb F_3$ is $-1$, so its order
is $26$: it does not divide $13$, and irreducibility excludes orders
$1,2$. Thus $E_L$ maps onto $k^\times$. The full global-unit image
modulo $k^\times$ has order $14$, as in the conductor calculation
above. Consequently the full image $U$ of $\mathcal O_M^\times$
in $K_{\rm res}^\times$ has order $364$ and contains $k^\times$.

Counting ordered bases gives
\[
 \#\operatorname{Gr}(3,6)(\mathbb F_3)
 =\frac{(3^6-1)(3^6-3)(3^6-9)}{(3^3-1)(3^3-3)(3^3-9)}=33880.
\]
The common saturation and scalar-isomorphism argument in
Theorem~\ref{pp:thm:fermatallquotients} applies to all these subspaces,
without imposing isotropy. Unpolarized classes are exactly $U$-orbits.
The stabilizer is $k^\times$ for the $28$ $k$-lines and
$\mathbb F_3^\times$ otherwise. Hence the orbit sizes are $14$ and
$182$, giving $28/14=2$ classes with ring $R$ and
$(33880-28)/182=186$ with ring $S$.

Set $E_0=\{u\in E_L:u\bmod\mathfrak P\in\mathbb F_3^\times\}$.
Then $[E_L:E_0]=13$. This odd index preserves the signature image:
$u^{13}\in E_0$ has the same signs as $u$. The image therefore has
order $32$, whereas $E_0/E_0^2$ has order $64$ by its rank five and
torsion $\{\pm1\}$. Thus $[E_0^+:E_0^2]=2$.
Furthermore $S^\times=\mu_6E_0$. Indeed, for $u\in S^\times$,
$u/\bar u$ is a root of unity reducing to $1$, hence belongs to
$\mu_3$; dividing by a suitable element of $\mu_6$ makes $u$ real
without changing the scalar residue condition. It follows that
$N_{M/L}(S^\times)=E_0^2$.

For any polarization on an $S$ quotient, pullback through the scalar
isogeny from $A$ gives a polarization on $A$. Its Rosati involution on
$M$ is conjugation by Proposition~\ref{pp:prop:fermatminpolarization};
transport through the rational isogeny proves the same for the quotient.
If an $S$ variety has one principal form $\lambda_0$, every other
principal form is $\lambda_0 a$ with $a\in E_0^+$: symmetry is
complex conjugation on $M$, and positivity is total positivity at
its real places. Its full unmarked automorphism group is $S^\times$,
so the classes are $E_0^+/E_0^2$, exactly two. The two $R$ varieties
also have two principal classes each by
Proposition~\externalref{packet}{pp:prop:actualpolarizations}.

Any principal form on any quotient pulls back to a degree-$729$
polarization on $A$. Proposition~\ref{pp:prop:fermatminpolarization}
transports it to $\lambda_{\min}$ by a source automorphism. Thus
Theorem~\ref{pp:thm:fermatallquotients} exhausts all principal polarization classes
on all these quotients, giving $160/2=80$ underlying varieties and
$188-80=108$ absolute principal obstructions.

For any minimum map to a principally polarizable target, the same
source normalization puts its kernel in $A[\pi]$. Source units commute
with $\pi$, so its original kernel is also in $A[\pi]$. Conversely,
every subgroup whose quotient is principally polarizable gives such
a minimum map. Orbit sizes now give
$2\cdot14+78\cdot182=14224$ kernels; the other
$108\cdot182=19656$ subgroups have nonprincipalizable quotients.
For the fixed pullback, each of the two principal classes contributes
seven kernels. They cannot share a kernel, since the common rational
form determines the quotient polarization uniquely.
\end{proof}

These counts concern complex abelian varieties and origin-preserving
maps. They neither identify new points of the normalized McMullen
curve nor assert a Jacobian classification for the $S$ varieties.

\begin{corollary}[Galois orbits after forgetting polarization]
\label{pp:cor:underlyinggalois}
The $80$ underlying principalizable varieties of
Theorem~\ref{pp:thm:fermatunderlying} form eleven Galois orbits.
One consists of $B,C$; the ten smaller-order orbits have sizes
\[
2,4,6,6,6,6,12,12,12,12.
\]
Each orbit in Theorem~\ref{pp:thm:quotientgalois} maps onto exactly one
of these underlying orbits, with fibers of size two.
\end{corollary}
\begin{proof}
Write $U_d$ for the subgroup of order $d$ in
$K_{\rm res}^{\times}$. The subgroup $U_{28}/U_{14}$, the kernel of the projection
$\overline G\to C_6\times C_2$ in the normalized Galois group,
acts trivially on underlying varieties,
because $U_{28}\subset U_{364}$ and $U_{364}$ is the full global-unit
residue image in the proof of Theorem~\ref{pp:thm:fermatunderlying}.
The unit lift gives an actual lattice isomorphism.

This kernel acts freely on polarized classes. A fixed class would
give $uH=vH$ with $u\in U_{28}\setminus U_{14}$ and $v\in U_{14}$.
The scalar $v^{-1}u$ has order $4$ or $28$, so generates a field of
degree $2$ or $6$ over $\mathbb F_3$. The invariant subspace $H$
would be a vector space over that field, contradicting
$\dim_{\mathbb F_3}H=3$. Each forgetful fiber already has exactly
two polarized classes, hence is exactly one free kernel orbit.
Equivariance prevents distinct polarized Galois orbits from merging
and halves every orbit size, as claimed.
The human-proof certificate is
\texttt{principalizable\_galois\_corollary\_v164}.
\end{proof}

This corollary concerns the principalizable subfamily only. It does
not classify Galois orbits of the $108$ nonprincipalizable varieties
or assert descent to a field of moduli.


\begin{proposition}[Cyclotomic action on all underlying quotients]
\label{pp:prop:allunderlyinggalois}
On all $188$ underlying classes of
Theorem~\ref{pp:thm:fermatunderlying}, the Galois action factors faithfully
through $\operatorname{Gal}(M/\mathbb Q)\simeq C_6\times C_2$.
Its orbits are exactly the full semilinear orbits on the
three-dimensional $\mathbb F_3$-subspaces of $K_{\rm res}=\mathbb F_{729}$,
modulo the global-unit subgroup $U_{364}$. Every orbit size divides $12$.
The principalizable and nonprincipalizable subsets are invariant.
\end{proposition}
\begin{proof}
The rational source $A$ and intrinsic subgroup $A[\mathfrak P]$ are
Galois-invariant, as in Theorem~\ref{pp:thm:quotientgalois}.
For every subgroup $H$, base change identifies the conjugate quotient
with $A/\sigma H$, and the actual normalizing unit $u_\sigma$ identifies
this with $A/(u_\sigma\sigma H)$. This requires no isotropy.
The common-saturation argument of
Theorem~\ref{pp:thm:fermatunderlying} identifies all underlying classes
with $U_{364}$-orbits. Choice ambiguities lie in $U_{14}\subset U_{364}$.
The proof of Theorem~\ref{pp:thm:quotientgalois} establishes surjectivity
onto the group of semilinear maps $G/U_{14}$ before its isotropic
orbit calculation, so applies on this full domain.

The subgroup $U_{364}$ is normal in
$\Gamma=K_{\rm res}^{\times}\rtimes\operatorname{Gal}(K_{\rm res}/\mathbb F_3)$.
For a generator $g$ of $K_{\rm res}^{\times}$, the scalar $g^{13}$ lies
in $G$, since $(g^{13})^{28}=-1$. Together with
$U_{364}=\langle g^2\rangle$ it generates all scalars.
Since $G$ surjects onto the six field automorphisms,
$\langle U_{364},G\rangle=\Gamma$. The quotient
$\Gamma/U_{364}\simeq C_2\times C_6$ has order $12$.

If $\sigma$ fixes $M$, its normalized action has $j=0$ and
$\epsilon=1$, hence scalar in $U_{28}\subset U_{364}$.
The actual action therefore factors through $\operatorname{Gal}(M/\mathbb Q)$.
Corollary~\ref{pp:cor:underlyinggalois} supplies an orbit of size $12$;
its point stabilizer in this order-$12$ group is trivial, proving
faithfulness. Principalizability is preserved by conjugation and its
inverse. The human-proof certificate is
\texttt{all\_underlying\_galois\_v165}.
\end{proof}

This determines the acting group on the $108$ nonprincipalizable
classes, but does not enumerate their orbits or prove descent of
individual varieties. There are $94$ full-scalar orbits on subspaces:
one of size $28$ and $93$ of size $364$, by the multiplier-field
stabilizers above. Frobenius may merge these scalar orbits.

\section{Minimum forms on nonprincipalizable quotients}
\begin{theorem}[Minimum polarizations on the obstructed quotients]
\label{pp:thm:obstructedminimum}
Each of the $108$ nonprincipalizable varieties of
Theorem~\ref{pp:thm:fermatunderlying} has minimum polarization degree
$3^{12}=531441$, with type $(1,3,3,3,3,9)$. Each has exactly $26$
minimum polarization classes under its full origin-preserving
 automorphism group.
\end{theorem}
\begin{proof}
Write $O=\mathcal O_M$, let $\mathfrak Q=\mathfrak P\cap\mathcal O_L$,
and fix the positive minimum form $E$ on $O$, with trace parameter
$\kappa_Oa_0$. Proposition~\ref{pp:prop:fermatminpolarization} gives
$(a_0)=\mathfrak Q$, $N_{L/\mathbb Q}(a_0)=-27$,
$\mathfrak QO=\mathfrak P$, and $(3)=\mathfrak Q^2$ in $\mathcal O_L$.
For $Y=A/H$, let $\Lambda$ be the inverse image of $H$ in
$\mathfrak P^{-1}$. Then $[\Lambda:O]=27$ and
$\Lambda O=\mathfrak P^{-1}$.

Pulling any polarization back to $A$ and transporting its Rosati
involution shows that every polarization on $Y$ has trace parameter
$\kappa_Oa$ with $a\in L$. Integrality on $O\times O$ gives
$a\in\mathcal O_L$. Mixed integrality on $O\times\Lambda$ is
equivalent to $a\Lambda\subset O$ by trace duality and conjugation.
Taking $O$-saturations gives $a\in\mathfrak P\cap L=\mathfrak Q$;
conversely this condition suffices for mixed integrality.
Thus $a=a_0v$ with $v\in\mathcal O_L$, and positivity is precisely
total positivity of $v$.

On $D=\mathfrak P^{-1}/O$ define, for $c\in k=\mathbb F_{27}$,
\[
 b_c(x+O,y+O)=3E(vx,y)\pmod3,\qquad v\bmod\mathfrak Q=c.
\]
This is independent of the lift and representatives: $E(O,\mathfrak
P^{-1})\subset\mathbb Z$ and $\mathfrak Q\mathfrak P^{-1}\subset O$.
The forms are alternating, and $b_c$ is nondegenerate for $c\ne0$.
Put
\[
 T_H=\{c\in k:b_c|_{H\times H}=0\}.
\]
This is an $\mathbb F_3$-linear subspace. Integrality on
$\Lambda\times\Lambda$ is equivalent to $v\bmod\mathfrak Q\in T_H$.
For an integral positive form $E_v(x,y)=E(vx,y)$, its degree on $Y$ is
\[
 \deg(E_v|_\Lambda)
 =\frac{729N_{L/\mathbb Q}(v)^2}{[\Lambda:O]^2}
 =N_{L/\mathbb Q}(v)^2.
\]
This determinant calculation does not assume that $E$ itself is
integral on $\Lambda$.

By Theorem~\externalref{supplement}{thm:unitnorm}, the residue image of $E_L^+$ is the
norm of the order-$364$ group $U$ from
Theorem~\ref{pp:thm:fermatunderlying}. Conjugation on $K_{\rm res}$
is the $27$th power map, so this norm image has order
$364/\gcd(364,28)=13$, exactly the squares in $k^\times$.
Since $-1$ is a nonsquare, every nonzero pair $\{c,-c\}$ contains
such a residue. If $T_H\ne0$, linearity therefore gives a totally
positive unit $\varepsilon$ with residue in $T_H$. Then $E_\varepsilon$
is a principal form on $Y$. Conversely a principal form has a unit
parameter and hence $T_H\ne0$. For the $108$ obstructed varieties,
$T_H=0$, and all positive integral parameters are exactly the totally
positive elements of $\mathfrak Q$.

Write $\mathfrak Q=(\beta)$ with $N_{L/\mathbb Q}(\beta)=-27$.
For such a parameter $v=\beta w$, positivity gives $N(w)<0$.
The signed norm bound of Proposition~\ref{pp:prop:fermatminpolarization}
gives $|N(w)|\ge27$, so $N(v)\ge729$. Equality is attained by $v=3$.
At equality $(v)$ has norm $3^6$, hence is $\mathfrak Q^2=(3)$,
since $\mathfrak Q$ is the unique prime above $3$, of norm $27$.
Thus all minimum parameters are exactly $3E_L^+$, and the minimum
degree is $729^2$.

The full automorphism group is $S^\times$, with norm image $E_0^2$
in the notation of Theorem~\ref{pp:thm:fermatunderlying}. The equal
signature images of $E_0$ and $E_L$ give $[E_L^+:E_0^+]=13$.
The minimum classes are therefore $E_L^+/E_0^2$, of size $13\cdot2=26$.

Finally let $F=3E_\varepsilon$ be a minimum form. Modulo $3$ it
vanishes on $O\times\Lambda$ and factors through $H=\Lambda/O$,
since $3\Lambda\subset O$. The induced alternating form is
$b_{\varepsilon\bmod\mathfrak Q}|_H$, nonzero because $T_H=0$.
It has rank $2$ on the three-dimensional space $H$. Thus exactly one
of the six elementary divisors of $F$ is prime to $3$. Their product
is $3^6$, so that divisor is $1$ and the other five have positive
$3$-adic exponents summing to $6$. The type is $(1,3,3,3,3,9)$.
\end{proof}

The residue obstruction also gives a direct count. If $H$ is isotropic
for two independent coefficients $c,d\in k$, then it is a Lagrangian
for $b_c$ and stable under $d/c$. This ratio generates $k$ over
$\mathbb F_3$, so $H$ is a $k$-line. Each of the $28$ such lines is
isotropic for every $b_c$. The $13$ coefficient lines therefore have
disjoint non-$k$-line Lagrangian populations, giving
$28+13(1120-28)=14224$ principalizable kernels and $19656$ obstructed
ones. This finite incidence identity corroborates the full-unit orbit
classification; the polarization minima additionally require the
signed norm argument above.

\begin{proposition}[Normalized group and minimum polarized action]
\label{pp:prop:minimumpolarizedgalois}
The normalized group $\overline G$ of
Theorem~\ref{pp:thm:quotientgalois} is $D_8\times C_3$, with $D_8$
of order $8$. Its commutator is the central involution $U_{28}/U_{14}$,
and its abelianization is $C_6\times C_2$.

The $2808=108\cdot26$ minimum polarized classes on the obstructed
quotients identify with the $U_{14}$-orbits of the $19656$ subspaces
with $T_H=0$, using the form descended from $3\lambda_{\min}$.
Their actual Galois action is the restriction of this same normalized
group, so every orbit size divides $24$. The central involution acts
freely and preserves the underlying variety. Each $26$-element
forgetful fiber is therefore the union of $13$ free pairs.
\end{proposition}
\begin{proof}
The scalar subgroup of $G$ is $U_{56}$; its quotient by $U_{14}$
is cyclic of order $4$, generated by $a$. Choose $b\in\overline G$
above the Frobenius generator of $C_6$. Then $bab^{-1}=a^3=a^{-1}$.
The cube of a representative is $w\operatorname{Frob}^3$ with
multiplier $\epsilon\in\{1,-1\}$. Its square is the scalar
\[
 w^{28}=\epsilon c^{26}=-\epsilon\in U_{14},
\]
since $c^{27}=-c$. Thus $b^6=1$, and its projection forces its
order to be $6$. The subgroups generated by $a$ and $b$ intersect
trivially and generate $\overline G$. Now $b^2$ is central of order
$3$, whereas $a,b^3$ generate $D_8$. These commuting factors
intersect trivially. Their commutator is $\langle a^2\rangle
=U_{28}/U_{14}$, proving the group assertions.

By Theorem~\ref{pp:thm:obstructedminimum}, all minimum forms on an
obstructed lattice are $3E_\varepsilon$ with $\varepsilon\in E_L^+$.
Choose $u\in\mathcal O_M^\times$ with $u\bar u=\varepsilon^{-1}$
using Theorem~\externalref{supplement}{thm:unitnorm}. Multiplication by $u$ from
$\Lambda_{u^{-1}H}$ to $\Lambda_H$ pulls $3E_\varepsilon$ back to
$3E$. Conversely $3E$ is integral on every such lattice and is
minimum when $T_H=0$. A fixed-pullback polarized isomorphism between
two quotients is a global unit $v$ by their common
$\mathcal O_M$-saturation. Preserving $3E$ gives $v\bar v=1$;
Kronecker's theorem then gives $v\in\mu_{42}$, with residue in
$U_{14}$. This proves the claimed equivalence relation.

An order-$7$ scalar generates $\mathbb F_{729}$ over $\mathbb F_3$
and cannot preserve a three-dimensional subspace. Every root orbit
therefore has size $7$, giving $19656/7=2808$. The source units
normalizing $\lambda_{\min}$ also normalize $3\lambda_{\min}$;
quotients and descended forms commute with base change. Hence
Theorem~\ref{pp:thm:quotientgalois} supplies the actual action here.
Nonprincipalizability is invariant under conjugation.

If $u\in U_{28}\setminus U_{14}$ fixes a root orbit, then $uH=vH$
for some $v\in U_{14}$. The scalar $v^{-1}u$ has order $4$ or $28$,
and generates a field of degree $2$ or $6$ over $\mathbb F_3$.
Its invariant subspace cannot have dimension $3$. Thus the involution
acts freely. Since $U_{28}\subset U_{364}$, it preserves the
underlying quotient, proving the fiber assertion. The human-proof
certificate is \texttt{obstructed\_minimum\_galois\_v165}.
\end{proof}

The $13$ pairs do not exhaust the forgetful fiber as a single kernel
orbit. Distinct polarized Galois orbits can merge after forgetting
the polarization; the division-by-two orbit conclusion of
Corollary~\ref{pp:cor:underlyinggalois} does not extend to these
$26$-element fibers. No orbit list or descent statement follows
from this proposition alone.

\section{Galois action fields and orbit classification}
\begin{theorem}[The ray field of the polarized-class action]
\label{pp:thm:quotientrayfield}
Let $\mathfrak P$ be the unique prime of $M$ above $3$, and put
\[
 \mathcal E=M\left(\sqrt{\frac{3+\sqrt{21}}2}\right).
\]
Then $\mathcal E$ is the ray class field of $M$ of modulus $\mathfrak P$,
with $[\mathcal E:M]=2$ and relative discriminant and conductor
$\mathfrak P$. It is Galois over $\mathbb Q$, with group
$D_8\times C_3$, where $|D_8|=8$.

The kernel of the normalized action on the $160$ principal polarized
classes is exactly $\operatorname{Gal}(\overline{\mathbb Q}/\mathcal E)$.
Their fields of moduli are contained in $\mathcal E$, with degrees
equal to their orbit sizes; a class in a size-$24$ orbit has field of
moduli $\mathcal E$. The $2808$ obstructed minimum polarized classes
also have fields of moduli contained in $\mathcal E$.
These statements do not assert effective descent of polarized models.
For the complete finite family, including $A$, the common fixing fields
are $M$ for the $189$ underlying classes and $\mathcal E$ for the
$2969$ minimum polarized classes.
\end{theorem}
\begin{proof}
The quadratic subfields $\mathbb Q(\sqrt{-3})$ and
$\mathbb Q(\sqrt{-7})$ of $M$ give $F=\mathbb Q(\sqrt{21})\subset M$.
Put $\omega=(1+\sqrt{21})/2$ and $a=\omega+1$. Then
$\omega^2-\omega-5=0$, $\omega^2-4=a$ and $N_{F/\mathbb Q}(a)=-3$.
In $F$, $(a)$ is the prime above $3$. Since $3$ has order $6$
modulo $7$, the prime in $M$ has ramification index $2$ and residue
degree $6$. Its ramification index over $F$ is $1$, so
$v_{\mathfrak P}(a)=1$ and $(a)\mathcal O_M=\mathfrak P$.

Let $t$ satisfy $t^2-\omega t+1=0$. The integral order
$\mathcal O_M[t]$ has discriminant $a$, which is a unit away from
$\mathfrak P$, including above $2$. At $\mathfrak P$, its discriminant
valuation $1$ precludes a nontrivial local index. Thus
$\mathcal E=M(t)$ is quadratic over $M$ with discriminant
$\mathfrak P$. The quadratic conductor--discriminant formula gives
the same conductor, with no infinite component because $M$ is CM
\cite[Theorem V.3.27]{MilneCFT2020}.
For an explicit equation, $\omega=t+t^{-1}$ gives the minimal polynomial
\[
 t^4-t^3-3t^2-t+1=0
\]
over $\mathbb Q$: $F\subset\mathbb Q(t)$ and $a$ is nonsquare even
in $M$, so the degree is $4$.

The class number of $M$ is $1$; independently, this follows from
Masley's Main Theorem for $m=21$, with $\varphi(21)=12\le24$ and
$21$ absent from its exceptional table \cite[p.~179]{Masley1976}.
The full-unit residue image is $U_{364}$ by
Theorem~\ref{pp:thm:fermatunderlying}. The finite-modulus ray sequence
therefore gives
\[
 \operatorname{Cl}_{\mathfrak P}(M)
 =\mathbb F_{729}^{\times}/U_{364}\simeq C_2.
\]
Ray class reciprocity and existence identify the conductor-$\mathfrak P$
quadratic extension just constructed with the ray field
\cite[Theorems V.1.7, V.3.6 and Remark V.3.8]{MilneCFT2020}.

Restrict the normalized homomorphism of
Theorem~\ref{pp:thm:quotientgalois} to
$\operatorname{Gal}(\overline{\mathbb Q}/M)$. Its image is
$U_{28}/U_{14}\simeq C_2$, hence it is a continuous abelian character.
For a local unit at a finite place $v$, form the idele equal to that
unit at $v$ and $1$ elsewhere. Its reflex norm is an everywhere
integral adelic unit. The same Tate-lattice and cyclotomic-factor
argument used in Theorem~\ref{pp:thm:quotientgalois} makes the CM
comparison map an actual polarized isomorphism. Consequently the
normalized residue action is the reflex norm modulo $U_{14}$.

If $v\ne\mathfrak P$, the reflex norm has component $1$ at
$\mathfrak P$: automorphisms preserve the underlying rational prime,
and $\mathfrak P$ is the only place above $3$. At $\mathfrak P$
the residue is $r^{286}$, as computed in
Theorem~\ref{pp:thm:quotientgalois}. This is trivial on principal units
and maps the full residue group onto $U_{28}$ because
$\gcd(286,728)=26$. The quadratic character has conductor exactly
$\mathfrak P$, by the local unit-filtration criterion
\cite[Section I.1.9]{MilneCFT2020}; its fixed extension of $M$ is
therefore $\mathcal E$.

The quotient of the normalized action by $U_{28}/U_{14}$ is precisely
$\operatorname{Gal}(M/\mathbb Q)$, so its full kernel is already
contained in $\operatorname{Gal}(\overline{\mathbb Q}/M)$.
It equals $\operatorname{Gal}(\overline{\mathbb Q}/\mathcal E)$.
Surjectivity and Proposition~\ref{pp:prop:minimumpolarizedgalois}
give the claimed degree and group over $\mathbb Q$. Faithfulness
on the $160$ classes identifies their common fixing field exactly.
The $2808$-class action factors through the same normalized map.
Taking stabilizer fixed fields proves the fields-of-moduli assertions.
No descent datum is supplied by this stabilizer calculation.
Finally $A$ is defined over $\mathbb Q$ and has a unique minimum
polarized isomorphism class, by
Proposition~\ref{pp:prop:fermatminpolarization}. Adding these fixed
classes changes neither kernel. Proposition~\ref{pp:prop:allunderlyinggalois}
gives the underlying common field $M$; the faithful $160$-class
subset detects the full polarized kernel while the $2808$-class
action factors through it. This proves the complete-family assertion.
\end{proof}

\begin{theorem}[Complete Galois orbit census of the quotient family]
\label{pp:thm:completequotientorbits}
The $188$ underlying quotient classes have $22$ Galois orbits, with
the following sizes and multiplicities:
\[
\begin{array}{c|rrrr}
\text{orbit size}&2&4&6&12\\\hline
\text{number of orbits}&2&4&4&12.
\end{array}
\]
The $108$ obstructed underlying classes form three orbits of size $4$
and eight of size $12$. Their $2808$ minimum polarized classes form
three orbits of size $8$ and $116$ of size $24$.
Above each underlying size-$4$ orbit lie one polarized size-$8$ orbit
and four size-$24$ orbits; above each size-$12$ orbit lie thirteen
size-$24$ orbits.

Including $A$, the complete family has $23$ underlying orbits and
$131$ minimum polarized orbits. The latter consist of one orbit of
size $1$, two of size $4$, four of size $8$, four of size $12$, and
$120$ of size $24$.
\end{theorem}
\begin{proof}
Write $K=\mathbb F_{729}$, $k=\mathbb F_{27}$ and $F(x)=x^3$.
Let $X$ be the $33880$ three-dimensional $\mathbb F_3$-subspaces
of $K$. We first count invariant subspaces of $T=uF^j$. For
$d=\gcd(6,j)$, with $d=6$ if $j=0$, put
$n=N_{K/\mathbb F_{3^d}}(u)$. Scalar conjugation changes $u$ by
$v^{1-3^j}$, whose image is precisely the norm-one subgroup. Thus
the fixed-space count depends only on $n$, and each norm value
occurs $728/(3^d-1)$ times.

For $j=0$, invariance makes a three-plane a vector space over
$\mathbb F_3(u)$, so only degrees $1$ and $3$ contribute.
The scalars $\pm1$ fix all $33880$ planes, the other $24$ elements
of $k^\times$ fix the $28$ $k$-lines, and all others fix none.

For $j=1,5$, the operator is cyclic with minimal polynomial $X^6-n$.
Indeed, a shorter relation among its powers would contradict linear
independence of the six field automorphisms over $K$. For $n=1$,
the two primary cyclic blocks $(X-1)^3,(X+1)^3$ have four submodules
of total dimension $3$. For $n=-1$, the single irreducible quadratic
factor of $(X^2+1)^3$ permits only even dimensions.

For $j=2,4$, a norm value outside $\mathbb F_3$ forces invariance
under $\mathbb F_9$, again excluding dimension $3$. For $n=\pm1$,
a sign and scalar conjugation reduce to $F^j$. The nilpotent
operator $F^j-1$ has cube zero and two-dimensional kernel, hence
Jordan type $(3,3)$. Its invariant three-planes number
\[
 \frac{3^6-3^4}{3^3-3^2}+\binom{2}{1}_{\!3}=36+4=40.
\]
Here the cyclic submodules have length $3$; those of type $(2,1)$
contain the kernel and correspond to lines of the second kernel
modulo the kernel. Type $(1,1,1)$ is impossible.

For $j=3$, $T$ is $k$-linear with $T^2=n$. If $n=1$, its two
eigenspaces have dimension $3$ over $\mathbb F_3$, giving
\[
 \sum_{r=0}^3\binom{3}{r}_{\!3}\binom{3}{3-r}_{\!3}
 =1+169+169+1=340
\]
invariant three-planes. If $n=-1$, an $\mathbb F_9$-structure
excludes them. Each of the twelve nontrivial square values in
$k^\times$ has degree $3$ and gives two invariant three-planes,
its eigenspaces. The eigenvalues $r,-r$ are not Frobenius conjugate:
$r^{3^i-1}=-1$ is impossible in the order-$26$ group for $i=1,2$,
since the power image has odd order. The remaining twelve nonsquare
values generate a degree-$6$ field after taking a square root and
give no invariant three-plane.

Let $S=X/K^\times$, of size $94$. The coset form of Burnside's
lemma says that the number of $F^j$-fixed elements of $S$ is the
average of $|\operatorname{Fix}_X(uF^j)|$ over $u\in K^\times$.
The preceding counts give
\[
\begin{array}{c|rrrrrr}
j&0&1&2&3&4&5\\\hline
|\operatorname{Fix}_S(F^j)|&94&2&10&14&10&2.
\end{array}
\]
For example the three nontrivial averages are
$364\cdot4/728=2$, $91\cdot2\cdot40/728=10$, and
$28(340+12\cdot2)/728=14$.
Fixed-point inversion gives respectively $2,4,4,12$ Frobenius
orbits of lengths $1,2,3,6$. Every full-scalar orbit splits into two
$U_{364}$-orbits: its stabilizer is either $\mathbb F_3^\times$ or
$k^\times$, both contained in $U_{364}$. Thus these lengths double
in the $188$-class action of
Proposition~\ref{pp:prop:allunderlyinggalois}. Subtracting
Corollary~\ref{pp:cor:underlyinggalois} gives the obstructed
underlying distribution.

For the polarized count, use the actual similitude group
$G=\{uF^j:u^{28}=\epsilon c^{3^j-1},\ \epsilon=\pm1\}$,
with $c^{27}=-c$ and $|G|=336$.
The sums of fixed three-planes restricted to these maps are
\[
\begin{array}{c|rrrrrr}
j&0&1&2&3&4&5\\\hline
\displaystyle\sum_{uF^j\in G}|\operatorname{Fix}_X(uF^j)|
&67760&112&560&9520&560&112.
\end{array}
\]
To verify the restriction, at $j=0$ only $\pm1$ contribute.
For odd $j=1,5$, $u^{364}=-\epsilon$, so the $28$ maps with
$\epsilon=-1$ each fix four planes. For $j=2,4$, the four possible
values $n=u^{91}$ satisfy $n^4=\epsilon$, each with multiplicity
$7$; only $n=\pm1$ contribute, giving $2\cdot7\cdot40$.
At $j=3$, $n=u^{28}=-\epsilon$, and the $28$ involutions each
fix $340$ planes. Burnside therefore gives $234$ $G$-orbits on $X$.

We remove the principalizable planes by incidence. The thirteen
coefficient lines in $k$ have pullback substitution
$t\mapsto t^{3^{-j}}$, up to the irrelevant scalar $\epsilon$.
Forward incidence transport uses its inverse; the fixed lines,
stabilizers and orbit partition are the same.
They form one fixed line $\mathbb F_3$ and four orbits of length
$3$. A line in a nontrivial orbit has stabilizer
$G_0=\{uF^j\in G:j=0,3\}$, of order $112$.
Scalar multiplication identifies its Lagrangians with those for
the actual form and normalizes $G_0$: at $j=3$ it changes $u$ by
$v^{26}$, whose $28$th power is $1$.

On the $1120$ Lagrangians, the $j=0$ fixed sum for $G_0$ is
$2\cdot1120$. At $j=3$, the maps with $\epsilon=1$ square to
$-1$ and fix no three-plane. Each of the $28$ maps with
$\epsilon=-1$ is an anti-symplectic involution with complementary
Lagrangian eigenspaces. An invariant Lagrangian is a subspace of
one eigenspace together with its annihilator in the other, giving
$\sum_r\binom{3}{r}_3=28$ choices. Hence $G_0$ has
$(2240+28\cdot28)/112=27$ Lagrangian orbits.

Consider pairs $(H,\ell)$ with $\ell\subset T_H$ a coefficient
line. Their number of $G$-orbits is $11+4\cdot27=119$, using
Theorem~\ref{pp:thm:quotientgalois} for the fixed coefficient line.
A generic principalizable plane occurs once; a $k$-line occurs
thirteen times. The $28$ $k$-lines are transitive under $G$ and
under its normal index-$3$ subgroup $G_0$, since the number of
$G_0$-orbits divides both $3$ and $28$. The incidence over these
lines has five orbits, while the lines themselves have one.
Thus the principalizable planes have $119-5+1=115$ orbits, leaving
$234-115=119$ obstructed polarized orbits.

The central involution acts freely, so stabilizers in $D_8\times C_3$
have orders $1,2,3,6$ and orbit sizes $24,12,8,4$. Let $h$ generate
the central $C_3$. Choose a representative $g$ above $F^2$ with
$g^3=1$. For each $v\in U_7$, the map $v^{-1}g$ also has cube
$1$, since $v^{91}=1$, and fixes $40$ three-planes. Their seven
fixed sets are disjoint: an overlap would make a three-plane
invariant under an order-$7$ scalar. These $280$ planes form
exactly $40$ root orbits fixed by $h$.

For a fixed principalizable plane outside the $k$-lines, its
unique coefficient line must be fixed by $F^2$, so it is the
actual coefficient line $\mathbb F_3$. The $k$-lines are also
isotropic. Hence the principalizable fixed root orbits are exactly
the fixed classes in the $160$-class population. Its two size-$4$
orbits and one size-$8$ orbit give $16$ fixed classes. Thus $h$
fixes $40-16=24$ obstructed polarized classes. If $a_d$ counts
their orbits of size $d$, then
\[
\begin{split}
 a_4+a_8+a_{12}+a_{24}&=119,\\
 4a_4+8a_8+12a_{12}+24a_{24}&=2808,\\
 4a_4+8a_8&=24.
\end{split}
\]
Nonnegativity gives $a_4=6-2a_8\ge0$ and
$a_{12}=2a_8-6\ge0$, forcing $a_8=3$, $a_{24}=116$ and
$a_4=a_{12}=0$.

Each polarized orbit maps onto an underlying orbit. The $312$
classes over an underlying size-$12$ orbit therefore form thirteen
size-$24$ orbits. Over each size-$4$ orbit the equation
$8x+24y=104$ forces $x\equiv1\pmod3$. There are three such
underlying orbits and only three size-$8$ polarized orbits in total,
so each has $x=1,y=4$. Finally add the known principal orbits and
the fixed minimum class of $A$. This proves all counts.
\end{proof}

\begin{corollary}[Fields of moduli in the obstructed family]
\label{pp:cor:obstructedmodulifields}
Put $F_0=\mathbb Q(\sqrt{-3},\sqrt{-7})$ and
$E_0=F_0(\sqrt{(3+\sqrt{21})/2})$.
The obstructed underlying classes in size-$4$ orbits have field of
moduli $F_0$; those in size-$12$ orbits have field of moduli $M$.
The minimum polarized classes in size-$8$ orbits have field of moduli
$E_0$, and those in size-$24$ orbits have field of moduli $\mathcal E$.
In particular their action is faithful, with common fixing field
$\mathcal E$.

The field $E_0$ is the degree-$8$ dihedral splitting field of
$X^4-X^3-3X^2-X+1$. None of these field-of-moduli statements
asserts effective descent.
\end{corollary}
\begin{proof}
The unique order-$3$ subgroup of $\operatorname{Gal}(M/\mathbb Q)$
fixes $F_0$; it is the stabilizer in a size-$4$ orbit. A size-$12$
orbit has trivial stabilizer. Similarly the unique central $C_3$
in $\operatorname{Gal}(\mathcal E/\mathbb Q)$ is the stabilizer in
a size-$8$ orbit. It fixes $F_0$ and cannot negate the square root
of $a=(3+\sqrt{21})/2$, since its order is $3$. Thus it fixes $E_0$.
Since $\sqrt a\notin M$, the degree of $E_0$ over $\mathbb Q$ is $8$,
so this is the full fixed field. A size-$24$ orbit has trivial stabilizer. Apply Theorem~\ref{pp:thm:quotientrayfield}.

Write $a=(3+\sqrt{21})/2$ and $a'=(3-\sqrt{21})/2$, so $aa'=-3$.
The field $E_0$ contains both $\sqrt a$ and
$\sqrt{a'}=\sqrt{-3}/\sqrt a$, hence all four roots
$t,t^{-1},t',t'^{-1}$. Conversely their sums give $\sqrt{21}$,
their differences give $\sqrt a,\sqrt{a'}$, and the product of
those differences gives $\sqrt{-3}$. The splitting field therefore
contains $F_0$ and $\sqrt a$, and equals $E_0$.
\end{proof}

\begin{theorem}[Fields of moduli for the complete finite family]
\label{pp:thm:completequotientmoduli}
Choose $\sqrt{21}>0$, put $a=(3+\sqrt{21})/2$, and set
\[
 D=\mathbb Q\left(\frac{(1+\sqrt{21})/2+\sqrt a}{2}\right),
 \qquad J=\mathbb Q(\zeta_7+\zeta_7^{-1}),
 \qquad M^+=\mathbb Q(\zeta_{21}+\zeta_{21}^{-1}).
\]
Retain $F_0=\mathbb Q(\sqrt{-3},\sqrt{-7})$,
$E_0=F_0(\sqrt a)$ and $\mathcal E=M(\sqrt a)$.
The fields of moduli of the complete quotient family, including $A$,
are given by the following tables. Non-normal fields are specified up to
$\mathbb Q$-conjugacy in $\mathcal E$.
\[
\begin{array}{c|rrrrr}
\text{underlying orbit size}&1&2&4&6&12\\\hline
\text{number of orbits}&1&2&4&4&12\\
\text{field of moduli}&\mathbb Q&\mathbb Q(\sqrt{21})&F_0&M^+&M
\end{array}
\]
\[
\begin{array}{c|rrrrr}
\text{minimum polarized orbit size}&1&4&8&12&24\\\hline
\text{number of orbits}&1&2&4&4&120\\
\text{field of moduli}&\mathbb Q&D&E_0&DJ&\mathcal E
\end{array}
\]
Here $D$ has degree $4$ and signature $(2,1)$, while $DJ$ has degree
$12$ and signature $(6,3)$. With the chosen embedding,
$DJ=\mathcal E\cap\mathbb R$; neither $D$ nor $DJ$ is totally real.
These remain fields of moduli, with no assertion of effective descent.
\end{theorem}
\begin{proof}
The normalized action projects to $\operatorname{Gal}(M/\mathbb Q)$
by $(j,\epsilon)$, where $\zeta_7\mapsto\zeta_7^{3^j}$ and
$\zeta_3\mapsto\zeta_3^\epsilon$. In particular complex conjugation
projects to $(3,-1)$. A stabilizer avoiding the free central involution
in $D_8\times C_3$ has order $1,2,3$ or $6$. Its order-$2$ factor,
if present, is a reflection. The two reflection conjugacy classes are
distinguished by the multipliers $\epsilon=1$ and $-1$ at $j=3$:
the two lifts over either pair differ by the central involution and
are conjugate by an order-$4$ scalar.

Write $K_{\mathrm{res}}=\mathbb F_{729}$ and let
$\operatorname{Frob}(x)=x^3$ on $K_{\mathrm{res}}$.
A reflection of multiplier $1$ fixes no root orbit of any three-plane
$H\subset K_{\mathrm{res}}$. Indeed write it as $g=u\operatorname{Frob}^3$. Then
$g^2=u^{28}=-\epsilon$. If $gH=vH$ for $v\in U_{14}$, the operator
$v^{-1}g$ preserves $H$ and has square $-\epsilon$, because $v^{28}=1$.
For $\epsilon=1$, its irreducible polynomial $X^2+1$ would make $H$
a vector space over $\mathbb F_9$, impossible in dimension $3$ over
$\mathbb F_3$. Thus every reflection stabilizer is conjugate to complex
conjugation $c$ in $\operatorname{Gal}(\mathcal E/\mathbb Q)$.

It follows that the stabilizers of the principal polarized orbits of
sizes $4,8,12,24$ are, up to conjugacy,
$C_3\times\langle c\rangle$, $C_3$, $\langle c\rangle$, and $1$.
The obstructed polarized orbits have only sizes $8,24$, already covered
by Corollary~\ref{pp:cor:obstructedmodulifields}.

The fields $E_0$ and $J$ lie in $\mathcal E$ and have coprime degrees
$8,3$, so $\mathcal E=E_0J$. By the quartic irreducibility in
Theorem~\ref{pp:thm:quotientrayfield},
$D=\mathbb Q(\sqrt{21},\sqrt a)$ has degree $4$.
It is fixed by $c$, as is $J$, and $[DJ:\mathbb Q]=12$.
Consequently $\mathcal E^{\langle c\rangle}=DJ$ and
$\mathcal E^{C_3\times\langle c\rangle}=D$.
The quartic has two real roots over $\sqrt{21}>0$ and one complex
pair over its negative conjugate. Since $J$ is totally real and
linearly disjoint from $D$, the signatures follow.

For underlying principalizable classes, forgetting the two forms is
exactly quotienting by the free central involution $z$.
Thus a polarized stabilizer $H$ gives the underlying stabilizer
$H\langle z\rangle/\langle z\rangle$ in $\operatorname{Gal}(M/\mathbb Q)$.
The size-$6$ underlying orbits therefore have fixed field $M^+$;
the size-$2$ orbits have the field fixed by $C_3$ and complex
conjugation, namely $\mathbb Q(\sqrt{21})$.
The size-$4$ and size-$12$ fields are $F_0$ and $M$.
Add the obstructed underlying orbits from
Corollary~\ref{pp:cor:obstructedmodulifields} and the rational class
of $A$, with its unique minimum polarization class.
The multiplicities are those of Theorem~\ref{pp:thm:completequotientorbits}.
\end{proof}

\section{Real models and actual descent}
\begin{theorem}[Real models and pseudoreal underlying varieties]
\label{pp:thm:quotientrealmodels}
In the complete $189$-class underlying family, exactly $15$ complex
isomorphism classes admit a model over $\mathbb R$. Each admits a model
over $\mathbb R_{\mathrm{alg}}=\overline{\mathbb Q}\cap\mathbb R$.
Another $14$ classes are \emph{pseudoreal}: they are isomorphic to their
complex conjugates but have no model over $\mathbb R$.

Exactly $28$ of the $160$ principal polarized quotient classes admit
real models, comprising two $R$-order and $26$ $S$-order classes.
They lie over $14$ underlying classes, one of order $R$ and $13$ of
order $S$, with both principal classes real on each. The $14$ pseudoreal
classes also comprise one $R$-order and $13$ $S$-order classes.
Including $A$ with its minimum polarization gives $29$ real minimum
polarized classes. None of the $2808$ obstructed minimum polarized
classes has a real model.

These count complex isomorphism classes admitting models, not inequivalent
real forms. Polarizations mean polarization morphisms; no chosen line-bundle
linearization or descent to the smallest number field is asserted.
\end{theorem}
\begin{proof}
Write $c$ for the chosen complex conjugation. The literal Fermat factor
$A$ is defined over $\mathbb Q$, and the restricted Jacobian polarization
$\mu$ is rational. Every geometric polarization $\lambda$ on $A$ has
$\lambda=\mu\alpha$ with $\alpha\in\End^0(A)=M$. Rosati is CM conjugation,
so symmetry gives $\alpha\in M^+$. The actual Galois action on $M$ is
cyclotomic, hence $c(\lambda)=\lambda$. In particular the minimum
polarization $\lambda_{\min}$ is defined over $\mathbb R_{\mathrm{alg}}$.
This does not identify it with $\mu$ or assert that it is rational.

Let $Y=A/H$ be a principal quotient with pullback $\lambda_{\min}$.
If its polarized class is fixed by $c$, the exact root-orbit equivalence
in Theorem~\ref{pp:thm:fermatallquotients} gives $c(H)=vH$ for
$v\in\mu_{42}$. Multiplication by $v^{-1}$ induces a polarized group
isomorphism $f:Y^c\longrightarrow Y$. Its conjugate cocycle is
\[
 f\,c(f)=1,\qquad v^{-1}\overline{v}^{\,-1}=1.
\]
This supplies a descent datum for the finite quadratic extension
$\overline{\mathbb Q}/\mathbb R_{\mathrm{alg}}$.
Since $Y$ is projective, finite locally free descent is effective
\cite[Lemma~39.25.3, tag 0CCJ]{StacksDescent}.
The origin, group law and polarization morphism descend by full
faithfulness for fpqc descent of morphisms
\cite[Section~35.35, tag 02VV]{StacksDescent}.
Conversely, a real polarized model forces its class to be fixed by $c$.

The real-embedding counts in
Theorem~\ref{pp:thm:completequotientmoduli} give
$1+2\cdot2+4\cdot6=29$ conjugation-fixed underlying classes and the
same number of minimum polarized classes including $A$.
There are $28$ principal fixed classes, two of order $R$ and $26$ of
order $S$; the obstructed minimum orbits have no fixed classes.
The free central involution exchanges the two principal classes over
each principalizable underlying variety and commutes with $c$.
Thus the $28$ real principal classes lie over $14$ underlying classes.
Together with $A$ this gives $15$ real underlying classes. On each of
the other $14$ conjugation-fixed underlying classes, $c$ exchanges its
two principal classes.

It remains to exclude even unpolarized real models for these last
classes. Suppose $Y_0/\mathbb R$ were such a model and choose a complex
isogeny $\phi:A_{\mathbb C}\longrightarrow Y_{0,\mathbb C}$.
The rank-one $M$-module of rational homomorphisms gives
$c(\phi)=\phi u$ for $u\in M^\times$. Commutativity of $M$ then gives
\[
 c(\phi\alpha\phi^{-1})
   =\phi u\overline\alpha u^{-1}\phi^{-1}
   =\phi\overline\alpha\phi^{-1}.
\]
So conjugation on the full endomorphism algebra of any such real form
must be CM conjugation. Choose a real polarization $\mu_0$ on $Y_0$.
For any geometric polarization $\eta=\mu_0\beta$, Rosati symmetry
puts $\beta\in M^+$. Consequently $c(\eta)=\eta$, and descent of
morphisms makes every geometric polarization morphism real.
In particular each principal class would be fixed, contradicting their
exchange. These $14$ classes are therefore pseudoreal. The remaining
$160$ underlying classes are not conjugation-fixed and cannot admit real
models. This proves the asserted counts and distinctions.
\end{proof}

\begin{corollary}[Real forms and forgetting the minimum polarization]
\label{pp:cor:quotientrealforms}
The family in Theorem~\ref{pp:thm:quotientrealmodels} has exactly
$58$ origin-preserving real minimum polarized forms: $56$ principal
forms, comprising four of order $R$ and $52$ of order $S$, and two
forms of the minimum-polarized source. There are exactly $29$ underlying
real forms. Forgetting the minimum polarization has two elements in
every fiber. All forms admit models over $\mathbb R_{\mathrm{alg}}$,
and each real polarized automorphism group is $\{1,-1\}$.
\end{corollary}
\begin{proof}
Fix a real model and let $\tau$ be its semilinear descent involution.
The preceding theorem forces conjugation on the full endomorphism
algebra to be CM conjugation. Its action on the full geometric polarized
automorphism group $G$ is therefore inversion. This group is $\mu_6$
for every principal quotient and $\mu_{42}$ for the source.
Every alternative polarized descent datum is $a\tau$ with $a\in G$,
and $(a\tau)^2=1$. An isomorphism $h$ between $a\tau$ and $b\tau$
exists exactly when $b=h^2a$. Effective descent thus identifies the
real forms with $G/G^2$, which has two classes represented by $1$ and
$-1$. Fixed automorphisms satisfy $h=h^{-1}$ and hence are $\pm1$.
This argument applies over both $\overline{\mathbb Q}/\mathbb R_{\mathrm{alg}}$
and $\mathbb C/\mathbb R$; the full automorphism groups agree under
algebraically closed base extension.

For the underlying variety with full order $T$, replace $G$ by
$T^\times$. The cocycle quotient is
\[
 \frac{\{u\in T^\times:u\bar u=1\}}
 {\{h/\bar h:h\in T^\times\}}.
\]
The numerator consists of roots of unity. Put $E_L=\mathcal O_{M^+}^{\times}$ and
$U_S=S^\times\cap M^+$. The adopted identities
$R^\times=\mu_6E_L$ and $S^\times=\mu_6U_S$ give denominator
$\mu_3$, so each of the $14$ real quotient complex classes has two
underlying real forms. For $T=\mathcal O_M$, root squares give
$\mu_{21}$ in the denominator, while the unit
$\gamma=z^6(1-z^2)$, $z=\zeta_{42}$, has
$\gamma/\bar\gamma=-z^{14}\notin\mu_{21}$. The denominator is
therefore all $\mu_{42}$, leaving just one underlying real form of $A$.
This gives $2\cdot14+1=29$ underlying forms.

For a fixed geometric principal class, the forgetful map on twist sets
is the identity $\mu_6/\mu_3\longrightarrow\mu_6/\mu_3$.
Each underlying quotient complex class has two principal classes,
so each underlying real form supports two real principal classes.
For the source the map is
$\mu_{42}/\mu_{21}\longrightarrow\mu_{42}/\mu_{42}$, with two
classes above its unique underlying real form. This proves the
$58$-to-$29$ distribution. The classification concerns polarization
morphisms and origin-preserving isomorphisms, without translations,
real torsors or higher-degree polarization counts.
\end{proof}

\section{A common field of definition}
\begin{theorem}[Common field of definition and underlying descent obstruction]
\label{pp:thm:quotientdefinitionfield}
Retain $M=\mathbb Q(\zeta_{21})$ and
$\mathcal E=M(\sqrt{(3+\sqrt{21})/2})$.
Let $Y$ be any of the $188$ underlying quotient classes $A/H$ in
Theorem~\ref{pp:thm:fermatunderlying}. For every number field
$M\subseteq F\subset\overline{\mathbb Q}$, the class $Y$ has a
model over $F$ if and only if $\mathcal E\subseteq F$.
In particular none of these $188$ classes descends to its embedded
field of moduli, all of which lie in $M$. This negative assertion
does not include the rational source $A$.

There is a single twist $A^{\mathrm{tw}}/\mathcal E$ carrying the
minimum polarization and with constant $\pi$-torsion. Its $33880$
order-$27$ subgroup kernels and quotient maps are defined over
$\mathcal E$. Hence $\mathcal E$ is the smallest common field of
definition inside $\overline{\mathbb Q}$ for the $189$ underlying
classes and also for their $2969$ minimum polarized classes.
The common underlying field of moduli is $M$, strictly smaller.
\end{theorem}
\begin{proof}
The source is rational with all endomorphisms defined over $M$.
Writing $\lambda_{\min}=\mu\alpha$ using its rational restricted
Jacobian polarization gives $\alpha\in M^+$, so $\lambda_{\min}$
is defined over $M$. On
$D=A[\pi]\simeq\mathbb F_{729}$, actual Galois action over $M$
is therefore a continuous scalar character
\[
 q:G_M\longrightarrow U_{28}.
\]
Indeed it is $\mathcal O_M/\mathfrak P$-linear and preserves the
polarization pairing, since $G_M$ fixes $\mu_3$. The normalization
unit in Theorem~\ref{pp:thm:quotientgalois} may here be chosen to be
$1$. Thus $q$ modulo $U_{14}$ is the adopted normalized action, and
Theorem~\ref{pp:thm:quotientrayfield} gives
\[
 \ker(q\bmod U_{14})=G_{\mathcal E}.
\]

Suppose $Y_0/F$ represents $A/H$, with $F\supseteq M$, and choose the
resulting degree-$27$ geometric isogeny $\phi:A\longrightarrow Y_0$.
The rank-one $M$-module of rational homomorphisms gives
$\sigma(\phi)=\phi a_\sigma$, $a_\sigma\in M^\times$, for
$\sigma\in G_F$. Since $G_F$ fixes $M$, this is a character.
The map $\phi$ is defined over a finite extension, so its character
has finite image and therefore lies in $\mu_{42}$. This finite-order
argument, rather than an assumed integrality of a quasi-isogeny ratio,
is what makes $a_\sigma$ an automorphism. Taking kernels shows
$\sigma(H)=a_\sigma^{-1}H$, so $G_F$ fixes the root orbit of $H$.
If some $q_\sigma$ lay outside $U_{14}$, a scalar in
$U_{28}\setminus U_{14}$ would preserve $H$. Such a scalar has
order $4$ or $28$ and generates a field of degree $2$ or $6$ over
$\mathbb F_3$. It cannot preserve a three-dimensional subspace.
Consequently $G_F\subseteq G_{\mathcal E}$ and
$\mathcal E\subseteq F$.

For sufficiency, the character $q$ restricted to $G_{\mathcal E}$
takes values in $U_{14}$. Reduction identifies
$\mu_{14}\subset\mu_{42}$ with $U_{14}$, since its kernel on
$\mu_{42}$ is $\mu_3$. Let $v_\sigma\in\mu_{14}$ be the unique
lift of $q_\sigma$. The roots are defined over $M$, so $v$ is a
character and $v_\sigma^{-1}\sigma$ is a coherent descent datum.
It factors through a finite Galois extension of $\mathcal E$.
Projectivity and descent of morphisms make it effective as a polarized
abelian variety $A^{\mathrm{tw}}/\mathcal E$ \cite{StacksDescent}.
The twisted action on $D$ is trivial. Every subgroup $H$ and quotient
map is consequently defined over $\mathcal E$
\cite[Chap.~IV, Lemma~2.1]{MilneAV2008}. For isotropic $H$,
the quotient polarization with pullback $\lambda_{\min}$ descends.
For obstructed $H$, use the minimum form with pullback
$3\lambda_{\min}$; these represent all obstructed minimum classes
by the root-orbit description above. The quotient construction
commutes with base change, and the uniquely determined polarization
morphism descends. This supplies every claimed model over $\mathcal E$.

Finally any common field of definition for all underlying classes
contains their common field of moduli $M$. The necessity just proved
then forces it to contain $\mathcal E$. For all minimum polarized
classes, their common field of moduli already is $\mathcal E$.
The simultaneous construction proves minimality in both cases.
No smallest field of definition for an individual class over fields
not containing $M$ is specified.
\end{proof}

The individual existence of CM models over $\mathcal E$ also follows
from the classical descent theorem of Milne, with its correction and
addendum \cite{MilneModuli1972,MilneModuli1974}.
That theorem concerns the fully marked polarized quotient triple,
whose field of moduli here is $\mathcal E$, rather than the coarse
underlying field contained in $M$. Its additional purity hypothesis
holds: the quotient root group $\mu_6$ is a direct summand of
$\mu_{42}$. The polarization-independent shorthand in Milne's
Remark~2 does not remove the polarization requirement on the
isomorphisms defining this field. Our construction gives, in addition,
one explicit twist with constant $\pi$-torsion and all quotient maps
simultaneously over $\mathcal E$; the underlying descent obstruction
is proved separately above.

\begin{theorem}[A common real number field for the real forms]
\label{pp:thm:realdefinitionfield}
Let $c$ be the chosen complex conjugation and put
$\mathcal E_{\mathbb R}=\mathcal E^{\langle c\rangle}
=\mathcal E\cap\mathbb R$. All $58$ minimum polarized real forms
and all $29$ underlying real forms classified above have models over
$\mathcal E_{\mathbb R}$ inducing those real forms under the selected
embedding. This field has degree $12$ and signature $(6,3)$; it is
not totally real.

Among the $14$ underlying real quotient complex classes, exactly $12$
have field of moduli $M^+$. For each of these, and for each of its real
forms, $\mathcal E_{\mathbb R}$ is the smallest real algebraic field
of definition inside $\overline{\mathbb Q}\cap\mathbb R$.
It is also the smallest common real algebraic field of definition
for all $15$ underlying real-model complex classes, for all $29$
underlying real forms, and for all $58$ minimum polarized real forms.
\end{theorem}
\begin{proof}
Use the character $q$ and its unique lift
$v:G_{\mathcal E}\longrightarrow\mu_{14}$ from
Theorem~\ref{pp:thm:quotientdefinitionfield}.
On $D$, conjugation is semilinear for $x\mapsto x^{27}$, so
\[
 q(c\sigma c)=q(\sigma)^{27}=q(\sigma)^{-1},
 \qquad v(c\sigma c)=c(v(\sigma)).
\]
The last equality follows from injectivity of root reduction on
$\mu_{14}$. The sequence from $G_{\mathcal E}$ to
$G_{\mathcal E_{\mathbb R}}$ splits by the actual involution $c$.
Consequently $v(\sigma)^{-1}\sigma$ on $G_{\mathcal E}$ together
with $c$ gives a coherent finite descent datum on the source.
The minimum polarization is defined over
$M^+\subset\mathcal E_{\mathbb R}$ and is preserved by these roots,
so the datum is effective as a polarized abelian variety.
Its $D$-action is trivial over $\mathcal E$ and is the original
conjugation action over $\mathcal E_{\mathbb R}$.
Only stability of $D$ is needed; the endomorphism $\pi$ itself is
not asserted to be defined over $\mathcal E_{\mathbb R}$.

For a conjugation-fixed principal quotient class choose $H$ and
$r\in\mu_{42}$ with $c(H)=rH$. Replace the source datum at $c$
by $r^{-1}c$. Its square is $1$, and its cross-relation with
$G_{\mathcal E}$ is unchanged because $r$ is defined over
$\mathcal E$ and commutes with the lifted roots. It preserves
$H$ and the source polarization. Finite subgroup quotient descent
and uniqueness of the quotient polarization therefore give the
desired principally polarized model over $\mathcal E_{\mathbb R}$.
For the source itself use the preceding datum.

Twisting each model by $-1$ along
$\mathcal E/\mathcal E_{\mathbb R}$ changes its real datum
from $\tau$ to $-\tau$: under the selected real embedding this
quadratic extension becomes $\mathbb C/\mathbb R$.
The root-square classification above identifies these as exactly
the two polarized real forms. This supplies all $58$ forms, hence
all $29$ underlying forms by the proved forgetful surjection.

The four polarized size-$12$ orbits have $24$ conjugation-fixed
classes by signature $(6,3)$. The two conjugation-fixed principal
classes above each real underlying class therefore give $12$ real
underlying classes with field of moduli
$M^+$. The other two real quotient classes, one of each order,
have field of moduli $\mathbb Q(\sqrt{21})$.
If an algebraic field $F\subset\overline{\mathbb Q}\cap\mathbb R$
defines one of the former $12$ classes, it contains $M^+$.
The model descends to a number field $F_0\subseteq F$, which we
enlarge inside $F$ to contain $M^+$. Base change to $F_0M$ and
Theorem~\ref{pp:thm:quotientdefinitionfield} imply
$\mathcal E\subseteq F_0M\subseteq FM$. Since $M/M^+$ is imaginary quadratic,
$(FM)^c=F$, whence $\mathcal E_{\mathbb R}\subseteq F$.
The constructed models prove the individual and common minimality
assertions. No minimality among arbitrary nonreal fields or
transcendental real extensions is asserted.
\end{proof}

\section{Minimum degrees and the remaining cases}
\begin{theorem}[Minimum degrees of underlying fields of definition]
\label{pp:thm:definitiondegrees}
Let $Y$ be one of the $188$ underlying quotient classes, with embedded
field of moduli $k_Y\subseteq M$. Every number field $F$ defining $Y$
satisfies $2\mid[F:k_Y]$. For the $144$ classes with $k_Y=M$, the
minimum absolute degree of a field of definition is $24$, attained
by $\mathcal E$. For the $24$ classes with $k_Y=M^+$, it is $12$,
attained by suitable conjugates of $\mathcal E_{\mathbb R}$.
The latter assertion includes the $12$ pseudoreal classes in those
orbits, using nonreal embedded conjugates.

For the remaining $16$ classes with field of moduli $B$, the lower
bound is $8$; for the remaining four with field of moduli
$\mathbb Q(\sqrt{21})$, it is $4$. No equality is asserted in
these last two cases. The rational source $A$ has minimum degree $1$.
\end{theorem}
\begin{proof}
A model over $F$ forces $k_Y\subseteq F$. Its base change to $FM$
and Theorem~\ref{pp:thm:quotientdefinitionfield} give
$\mathcal E\subseteq FM$, so $[FM:M]$ is even.
Since $M/\mathbb Q$ is Galois, restriction identifies
$\operatorname{Gal}(FM/F)$ with
$\operatorname{Gal}(M/(M\cap F))$. The degree formula gives
\[
 [FM:M]=[F:M\cap F],\qquad
 [F:k_Y]=[F:M\cap F]\,[M\cap F:k_Y].
\]
Thus $[F:k_Y]$ is even. The fields-of-moduli table gives the lower
bounds and the counts $144+24+16+4=188$.

The common $\mathcal E$ models attain degree $24$ when $k_Y=M$.
Each of the four underlying size-$6$ orbits contains three real-model
classes: its polarized size-$12$ orbit has six conjugation-fixed
classes, with two above each real underlying class.
Theorem~\ref{pp:thm:realdefinitionfield} supplies an
$\mathcal E_{\mathbb R}$ model for such a real member. Conjugating
this model by an automorphism of $\overline{\mathbb Q}/\mathbb Q$
supplies a degree-$12$ model for every member of the same orbit.
For a pseudoreal member the resulting embedded field cannot lie in
$\mathbb R$, although its abstract signature is still $(6,3)$.
This proves degree minimality without asserting a unique smallest
field under inclusion among arbitrary nonreal fields.
\end{proof}

\section{Weil structures and maps from obstructed quotients}
\begin{corollary}[Weil structures and Hodge classes on the quotient family]
\label{pp:cor:quotientweilhodge}
All $160$ principal polarized classes of
Theorem~\ref{pp:thm:fermatallquotients}, every minimum polarization
on each of the $108$ obstructed varieties of
Theorem~\ref{pp:thm:obstructedminimum}, and the minimum class on the
actual Fermat factor $A$ have integral
$\mathcal O_K$ action for $K=\mathbb Q(\sqrt{-3})$, signature $(3,3)$,
and Hermitian discriminant class $-1$.

Each of the $188$ underlying quotients, and $A$, has rational Hodge
algebra generated by divisors and its two-dimensional $K$-Weil space
in codimension three, with Hodge-dimension polynomial
$(1+t)^6+2t^3$. Every rational Hodge class on every finite product of
these varieties, with repetitions allowed, is algebraic.
\end{corollary}
\begin{proof}
Choose the minimum form $E$ on $A$ as the pullback of an actual
principal B form through a degree-$27$ $M$-linear isogeny of
Proposition~\ref{pp:prop:fermatbcminimum}. The accepted B form has
signature $(3,3)$ and discriminant $-1$ by
Corollary~\externalref{packet}{pp:cor:weilimage}. Pullback gives a rational $K$-linear
isometry, so $E$ has the same discriminant. All $160$ principal
classes admit a quotient presentation with pullback $E$ by
Theorem~\ref{pp:thm:fermatallquotients}; the required source unit
normalization commutes with $K$. Their rational Hermitian spaces
are therefore isometric to B's.

The full integral orders $O,R,S$ all contain
$\mathcal O_K=\mathbb Z[z^{14}]$, since $z^{14}\equiv1\pmod{\mathfrak P}$.
The actual analytic $K$ action retains signature $(3,3)$ under the
specified isogenies, and every polarization Rosati involution restricts
to conjugation. For a minimum form $F=3E_\varepsilon$ on an obstructed
variety, Theorem~\ref{pp:thm:obstructedminimum} gives
$\varepsilon\in E_L^+$. In common rational coordinates its Hermitian
form is $h_F(x,y)=h_E(3\varepsilon x,y)$. Its determinant ratio is
\[
 N_{M/K}(3\varepsilon)=N_{L/\mathbb Q}(3\varepsilon)
 =3^6=729=N_{K/\mathbb Q}(27),
\]
because $M=LK$, $L\cap K=\mathbb Q$, and $N(\varepsilon)=1$.
The discriminant class remains $-1$. The minimum forms on $A$ differ
by positive real units and have determinant ratio $1$. None of these
rational determinant statements changes the nonprincipal polarization
types or the integral lattices.

Every quotient is $M$-linearly isogenous to $A$ and hence B. These
algebraic isogenies identify rational cohomology rings, divisor subspaces,
and the rational Weil spaces
$W_K(Y)=\bigwedge_K^6H^1(Y,\mathbb Q)$. Thus the B-type
 divisor/Weil generation and the displayed dimension polynomial transfer
from the all-powers argument of Section~\externalref{hodge}{ext:sec:main}.
The principal B model and the original trace-polarized B fiber are
isogenous, as established in the principal-lattice construction; their
integral polarizations are not being identified.

A product $X$ of $n$ chosen factors is isogenous to $B^n$. Taking
rational quasi-inverses in $\operatorname{Hom}\otimes\mathbb Q$ and
clearing denominators gives an actual isogeny $f:B^n\to X$.
Theorem~\externalref{hodge}{ext:thm:individual} supplies rational algebraicity on
$B^n$. If a rational algebraic cycle $Z$ represents $f^*\alpha$,
then $(\deg f)^{-1}f_*Z$ represents $\alpha$, since
$f_*f^*=\deg f$ on rational cohomology. This proves algebraicity in
every codimension for the stated products, all in one isogeny class.
\end{proof}

This corollary supplies rational algebraicity and ambient Weil structures.
It gives no additional McMullen curve points, identification with
Mostaed's literal normalization, integral Hodge statement, or explicit
cycle equations.

\begin{theorem}[Minimum principal-target isogenies from obstructed quotients]
\label{pp:thm:obstructedisogenies}
For each of the $108$ obstructed varieties $Y=A/H$, the minimum degree
of an isogeny to any principally polarized target is $729$.
The minimum degree of $Y\to A$ is $27$; all such minimum maps have
one kernel $C_0\simeq(\mathbb Z/3)^3$ and differ by target automorphisms.

For any fixed minimum polarization $\lambda$ on $Y$, the maps to
principal targets with pullback $\lambda$ have exactly $89{,}360{,}320$
distinct kernels. Their group types and counts are
\[
\begin{array}{c|r}
(\mathbb Z/3)^6&91{,}840\\
\mathbb Z/9\oplus(\mathbb Z/3)^4&89{,}268{,}480.
\end{array}
\]
Exactly $1120$ contain $C_0$, equivalently factor through a minimum
map $Y\to A$; the remaining $89{,}359{,}200$ do not.
\end{theorem}
\begin{proof}
Use the common lattices $O$ and $\Lambda$ for $A$ and $Y$, with
$[\Lambda:O]=27$ and $\Lambda O=\mathfrak P^{-1}$.
Every rational map is scalar in $M$. The condition $c\Lambda\subset O$
is equivalent, by taking $O$-saturations, to $c\in\mathfrak P$.
Conversely $\mathfrak P\Lambda\subset O$, so
\[
 \operatorname{Hom}(Y,A)=\mathfrak P=\pi O,\qquad
 \deg(\pi u:Y\to A)=27|N_{M/\mathbb Q}(u)|.
\]
The minimum is $27$, attained precisely for $u\in O^\times$.
These maps differ by target units and have kernel
$C_0=\mathfrak P^{-1}/\Lambda$, of order $27$ and exponent $3$.
This is the kernel of $\pi:Y\to A$, not of $\pi:Y\to Y$.

A degree-$d$ map to a principal target pulls back a polarization
of degree $d^2$. Theorem~\ref{pp:thm:obstructedminimum} gives
$d\ge729$. Composing $\pi:Y\to A$ with a degree-$27$ principal
quotient of $A$ attains this bound. More precisely, write the fixed
minimum form on $Y$ as $F=3E_\varepsilon$, with
$\varepsilon\in E_L^+$. Since $\pi\bar\pi=3$, the map $\pi$
pulls $E_\varepsilon$ back to $F$. Thus the bound is attained with
each prescribed minimum form as pullback.

Let $D=\Lambda^\#/\Lambda$ be the discriminant group for $F$.
Its elementary divisors give the symplectic module
\[
 D\simeq(\mathbb Z/9)^2\oplus(\mathbb Z/3)^8,
\]
with one hyperbolic pair of denominator $9$ and four of denominator
$3$. Integral self-dual overlattices correspond precisely to isotropic
subgroups $J=J^\perp$ of order $3^6$, by the dual-lattice argument of
Theorem~\ref{pp:thm:fermatallquotients}. These are exactly the kernels
with fixed pullback $F$, including all possible principal targets.

The subgroup $D[3]$ has dimension $10$ over $\mathbb F_3$, with
pairing radical $3D$ of dimension $2$. An elementary Lagrangian of
dimension $6$ must contain this radical: its image in the symplectic
quotient $D[3]/3D$ has dimension at most $4$. Such Lagrangians are
therefore the inverse images of the Lagrangians in that eight-dimensional
quotient, giving
\[
 L_4(3)=(3+1)(9+1)(27+1)(81+1)=91840.
\]

For a cyclic subgroup $C=\langle x\rangle$ of order $9$,
nondegeneracy supplies $y$ with pairing $1/9$. The pair $x,y$
generates a nondegenerate $(\mathbb Z/9)^2$ summand $V$, so
$D=V\oplus W$ orthogonally. Since $3V$ already has order $9=|3D|$,
we have $3W=0$; hence $W$ is symplectic $\mathbb F_3^8$.
Moreover $C^\perp=C\oplus W$. The Lagrangians containing $C$
are exactly $C\oplus L$, with $L$ a Lagrangian in $W$, so there
are $91840$ and each has type $\mathbb Z/9\oplus(\mathbb Z/3)^4$.
Every non-elementary Lagrangian contains an order-$9$ element,
so this exhausts the remaining types.

There are
\[
 \frac{3^{12}-3^{10}}{9-3}=78732
\]
cyclic subgroups of order $9$ in $D$. Each non-elementary Lagrangian
contains $(3^6-3^5)/(9-3)=81$ such subgroups. Counting pairs
$(C,J)$ with $C\subset J$ therefore gives
\[
 \frac{78732\cdot91840}{81}=89268480
\]
non-elementary kernels. Adding the elementary population gives the
stated total.

Finally, a fixed-form principal quotient factors through
$r=\pi:Y\to A$ exactly when its kernel contains $C_0$.
The resulting map from $A$ has degree $27$ and pullback
$E_\varepsilon$, since pullback by $r$ is injective on rational forms.
Conversely every such principal quotient of $A$ gives a kernel above
$C_0$. The subgroup correspondence $J\mapsto r^{-1}(J)$ is bijective,
so Theorem~\ref{pp:thm:fermatallquotients} gives exactly $1120$.
All minimum maps $Y\to A$ differ by target automorphisms, hence give
the same factorization condition.
\end{proof}

The $89{,}360{,}320$ count fixes a minimum pullback form. It counts
subgroups, not polarized target classes, and is not multiplied by
$26$ when the pullback form varies: intersections between those
kernel populations require a separate analysis.

The finite family consisting of $A$ and its $188$ underlying quotients
has the following complete minimum-polarization data. Here $c_{\min}$
counts minimum polarization classes on one underlying variety, and
$d_{\rm pp}$ is the minimum isogeny degree to a principal target.
\begin{center}
\small
\begin{tabular}{c|r|r|c|r|r}
Full order & $\#Y$ & $\deg\lambda_{\min}$ & Type & $c_{\min}$ & $d_{\rm pp}$\\\hline
$O$ & 1 & 729 & $(1,1,1,3,3,3)$ & 1 & 27\\
$R$ & 2 & 1 & $(1,1,1,1,1,1)$ & 2 & 1\\
$S$ & 78 & 1 & $(1,1,1,1,1,1)$ & 2 & 1\\
$S$ & 108 & 531441 & $(1,3,3,3,3,9)$ & 26 & 729
\end{tabular}
\end{center}
The entries follow from Proposition~\ref{pp:prop:fermatminpolarization}
and Theorems~\ref{pp:thm:fermatunderlying},
\ref{pp:thm:obstructedminimum} and~\ref{pp:thm:obstructedisogenies}.
Thus the $189$ underlying varieties support exactly
$1+2\cdot2+78\cdot2+108\cdot26=2969$ minimum polarized classes.
Different underlying varieties cannot merge as polarized varieties.
Corollary~\ref{pp:cor:quotientweilhodge} gives their integral Weil
structures and rational Hodge algebraicity on finite products.
This is a finite Fermat three-torsion quotient family, distinct from
the corrected $\ell=43$ Hecke packet and from a curve-point census.

\section{A structural two-prime comparison}
\begin{proposition}[The two-prime primitive factor]
\label{pp:prop:twoprimefermat}
Let $p<q$ be distinct odd primes, $m=pq$, and let $H_d$ denote
the smooth projective curve $y^2=x^d-1$. In $J(H_m)$ let $A_m$ be
the connected image of the primitive $\Phi_m$ projector, with its
actual inclusion. Its dimension is $(p-1)(q-1)/2$, and the canonical
restricted polarization has type
\[
 \bigl(1^{[(p-2)(q-1)/2]},\;
       p^{[(q-p)/2]},\;(pq)^{[(p-1)/2]}\bigr),
\]
where bracketed superscripts denote multiplicities.
Both this polarization and the addition isogeny
$A_m\times J(H_p)\times J(H_q)\to J(H_m)$, using pullback along
$(x,y)\mapsto(x^q,y)$ and $(x,y)\mapsto(x^p,y)$, have degree
$q^{p-1}p^{q-1}$.
For the covariant rotation $T$ of $J(H_m)$, put
\[
 \mathcal N_m=pq-p\sum_{j=0}^{q-1}T^{pj}
                  -q\sum_{j=0}^{p-1}T^{qj}.
\]
Then $\operatorname{im}\mathcal N_m=A_m$, $\mathcal N_m$ is
Rosati symmetric, and $\mathcal N_m^2=pq\mathcal N_m$.
The least positive denominator of the projector
$e_m=\mathcal N_m/(pq)$ in $\operatorname{End}(J(H_m))$ is $pq$.
\end{proposition}
\begin{proof}
The prime-degree ramification argument of
Proposition~\ref{pp:prop:fermatrestricted} applies to both actual
pullbacks: they are embeddings with induced forms $q\lambda_{H_p}$
and $p\lambda_{H_q}$. Their integral separation holds uniformly.
Indeed choose $1\leq a<p$ with $qa=1+kp$. The identity
\[
 \left(\sum_{j=0}^{a-1}t^{qj}\right)\Phi_q(t)
 -t\left(\sum_{l=0}^{k-1}t^{pl}\right)\Phi_p(t)=1
\]
follows by expanding the two consecutive geometric sums.
Each cyclotomic polynomial annihilates the corresponding lower
Jacobian image, so those images have zero intersection, including
torsion. Their character spaces are orthogonal. The lower product's
polarization kernel is consequently
\[
 (\mathbb Z/q)^{p-1}\oplus(\mathbb Z/p)^{q-1}
 \simeq(\mathbb Z/p)^{q-p}\oplus(\mathbb Z/pq)^{p-1}.
\]
The holomorphic differentials $x^{i-1}dx/y$ and their conjugates
give every nontrivial character of the order-$m$ rotation group once in full
cohomology. The remaining characters have order $pq$, so the actual
$A_m$ is the connected orthogonal complement of the lower product.
The saturated unimodular lattice argument of
Proposition~\ref{pp:prop:fermatrestricted} transfers the displayed
discriminant group to $A_m$. Its nonunit elementary divisors give
$(q-1)/2$ entries; the dimension supplies the stated number of unit
entries. The common group order gives both claimed degrees.

The two deck sums act by $q$ and $p$ on their respective lower
character spaces and vanish on the primitive part. Thus
$e_m$ is the orthogonal primitive projector. Inversion preserves
the deck sums, proving Rosati symmetry. Its polynomial residual is
$\mathcal N_m^2-pq\mathcal N_m=2pq(1+T+\cdots+T^{pq-1})$;
the full trace factors through the Jacobian of the genus-zero rotation
quotient and therefore vanishes. Finally, the discriminant group has exponent $pq$.
The lattice-projection proof of
Corollary~\ref{pp:cor:fermatprojector} therefore proves that $pq$
is the least integral denominator.
\end{proof}

This uniform statement concerns the induced polarization and the
specified projector. It requires neither simplicity of $A_m$ nor
a classification of its other polarizations.



% !TeX root = ../../quotients-and-descent.tex
\section{The restricted principal census and Jacobian question}\label{q:jacobians}
This section supplies the order-and-type result invoked for the source
factor. Its arithmetic proof is independent of the quotient and descent
sections. The theta step has an additional execution and enclosure
contract; the exact conditional statement is retained below.

\begin{theorem}[Complete simple principal classification with the selected order]
\label{pp:thm:tenprincipal}
There are exactly ten origin-preserving polarized isomorphism classes of
simple principally polarized complex abelian sixfolds $X$ such that
$\operatorname{End}^0(X)\simeq M$ and, under this identification,
$R\subseteq\operatorname{End}(X)$. The containment condition is independent
of the field identification. Eight classes have endomorphism ring $R$
and polarized automorphism group $\mu_6$; two have endomorphism ring
$\OM$ and polarized automorphism group $\mu_{42}$. They belong to four
unpolarized isogeny classes, with the following counts at one CM-type
representative from each orbit:
\[
\begin{array}{c|c|ccc}
\Phi&\text{flip parity from }\Phi_0 &[R]&[\mathfrak a]&[\OM]\\ \hline
\{1,5,11,13,17,19\}&0&2&2&0\\
\{1,5,11,13,19,25\}&1&0&0&1\\
\{1,5,11,17,23,29\}&0&2&2&0\\
\{1,5,11,23,25,29\}&1&0&0&1
\end{array}
\]
In particular, the four specified polarized B/C objects exhaust the
objects in their unpolarized isogeny class satisfying this $R$-containment
condition. No assertion is made here for other endomorphism orders or
for the induced A type.
\end{theorem}
\begin{proof}
Every automorphism of $M$ preserves $\OL$ and takes $1-z^{14}$ to
itself or an associate. Hence it preserves $R=\OL+(1-z^{14})\OM$.
After choosing a field identification, $H_1(X,\Ql)$ has dimension one
over $M$, and its integral lattice is a full fractional $R$-ideal.
Conversely, the positive unimodular forms of
Proposition~\externalref{packet}{pp:prop:markedpolarizations} give abelian varieties
by the Riemann criterion.

We check full endomorphisms and simplicity, not merely a centralizer.
For a CM type $\Phi$, put
$H=\{\sigma\in\operatorname{Gal}(M/\Ql):\sigma\Phi=\Phi\}$.
The splitting of
$\operatorname{End}_{\Ql}(M)=\bigoplus_\sigma M\sigma$
over $\mathbb C$ has disjoint permutation-matrix supports for distinct
$\sigma$. A nonzero coefficient in $M$ has nonzero value at every
embedding, so a term crossing the two Hodge blocks cannot cancel.
Thus the full rational Hodge endomorphism algebra is
$M\mathbin{\rtimes}H=\operatorname{End}_{M^H}(M)$.
The rational weight-one equivalence
\cite[Corollary~6.9]{MilneShimura2017} and Poincar\'e complete
reducibility show that the abelian variety is simple precisely when
$H$ is trivial: otherwise this matrix algebra has nontrivial
idempotents, whereas for trivial $H$ it is the field $M$.

The exact census of the $64$ types gives precisely the four
trivial-stabilizer orbits in the table, each of size $12$.
The two retained implementations are
\path{experiments/fixed_field_cm_census_v98/primary.py} and
\path{experiments/fixed_field_cm_census_v98/independent.py}.
The bounded table derivation
\path{experiments/simple_cm_unmarked_v145/derive.py}
binds their accepted outputs, checks orbit coverage and parity, and
applies Proposition~\externalref{packet}{pp:prop:markedpolarizations}; it does not
re-prove the ideal classification or the unit-norm theorem.

Every polarization on a simple object here has Rosati involution
equal to conjugation on $M$. Indeed its positive trace form extends
to $M\otimes\mathbb R\simeq\mathbb C^6$. An involution exchanging
factors would give trace norm zero on a factor idempotent. On a fixed
complex factor, the identity would give negative trace norm on $i$.
Positivity therefore forces conjugation on every factor. Thus the
marked classification covers every principal polarization.

Fix one primitive type representative. An unmarked polarized
isomorphism conjugates the full endomorphism fields and hence is
semilinear for some $\sigma\in\operatorname{Gal}(M/\Ql)$.
Holomorphicity at this fixed type forces $\sigma\Phi=\Phi$,
so $\sigma=1$. Forgetting the marking therefore identifies no
additional classes or automorphisms at the chosen representative.
Every object in its orbit can be remarked to this representative,
because $R$ is Galois-stable. The table and automorphism assertions
now follow from Proposition~\externalref{packet}{pp:prop:markedpolarizations}.

An unpolarized isogeny likewise induces a field automorphism relating
the types, so distinct primitive orbits are not isogenous. Within one
orbit, after remarking, full fractional lattices in $M$ are
commensurable; an integer multiple of the identity gives an isogeny.
Finally $\sigma_{31}(\Phi_C)=\Phi_B$ places the four distinct
objects of Proposition~\externalref{packet}{pp:prop:actualpolarizations} in the first
orbit. Its count of four proves the restricted exhaustion assertion.
\end{proof}

\begin{corollary}[Jacobian membership in the ten-class list]
\label{pp:cor:jacobianlist}
Of the ten classes in Theorem~\ref{pp:thm:tenprincipal}, the four in
the third row and the single class in the second row are not
canonically polarized Jacobians of smooth curves. The single class
in the fourth row is the known Jacobian of the smooth projective
model of $y^7=x^3-1$.
These signature and automorphism criteria alone leave the four B/C
classes in the first row undecided; Theorem~\ref{pp:thm:jacobiancomplete}
will exclude them using the invariant theta divisor.
If any is a Jacobian, its curve is nonhyperelliptic, has full
automorphism group $C_3$, and is a cyclic triple cover of
$\mathbb P^1$ with eight branch points, four of each nonzero
inertia exponent.
\end{corollary}
\begin{proof}
Strong Torelli \cite[Theorem~12.1]{MilneJacobians2021} realizes an
odd-order polarized Jacobian automorphism $\beta$ by a curve
automorphism up to sign. Squaring the latter induces $\beta^2$
and has the same odd order, by faithfulness on the Jacobian.
For $\beta=\zeta_3$, there are no invariant differentials, so the
curve quotient has genus zero. Riemann--Hurwitz gives eight fixed
points. If $n_1,n_2$ count the two inertia exponents, then
$n_1+n_2=8$ and $n_1+2n_2\equiv0\pmod3$.
The cyclic-cover eigenspace formula
\cite[equation~(2.2)]{LiMantovanPries2021} gives only signatures
$(4,2),(3,3),(2,4)$, respectively for $n_2=1,4,7$.
The third row has signature $(1,5)$ and is therefore excluded.
This is also directly Zarhin's established criterion
$g+2<3|a-b|$ \cite[Theorem~1.1]{Zarhin2007}, here $8<12$;
the obstruction method is not new.
The first row has signature $(3,3)$ and hence requires $n_1=n_2=4$.
A hyperelliptic curve would induce a nontrivial order-three map on
its hyperelliptic base $\mathbb P^1$, with only two fixed points
and at most four curve points above them. This contradicts eight
fixed points. Strong Torelli and the full polarized automorphism
group $\mu_6$ then give full curve automorphism group $C_3$.

For either maximal-order object, use instead $\beta=\zeta_{21}$.
Its eigencharacters are primitive, so again the curve quotient
has genus zero. If $a,b,c$ count branch points with ramification
indices $3,7,21$, Riemann--Hurwitz gives
$7a+9b+10c=26$. Its only nonnegative solution is $(1,1,1)$.
After fixing these three branch points at $0,1,\infty$, the
monodromy residues are $(7u,3v,-7u-3v)$, with
$u\in(\mathbb Z/3)^\times$ and $v\in(\mathbb Z/7)^\times$.
The Chinese remainder theorem makes these twelve tuples a single
orbit under a change of cyclic generator. Kummer theory thus gives
one curve up to complex isomorphism, namely $y^7=x^3-1$.
Its known simple maximal-order CM type is
$\{1,2,4,8,10,16\}$ modulo $21$
\cite[Proposition~6.11, Lemma~6.12 and Proposition~6.14]{LiMantovanPries2021}.
Its odd lifts modulo $42$ are $\{1,23,25,29,31,37\}$, which are
$23$ times the fourth row's representative. That row is therefore
the known Jacobian. The distinct second-row object cannot also be
a Jacobian. The finite ramification and type arithmetic is retained
in \path{experiments/simple_cm_unmarked_v145/jacobian_check.py};
the proof uses the geometric inputs just named and makes no
assertion about isogeny to a Jacobian.
\end{proof}

\subsection{An invariant theta obstruction and the Jacobian classification}
\label{pp:sec:thetacertificate}
Write $B_0,B_1,C_0,C_1$ for the four principal objects of
Proposition~\externalref{packet}{pp:prop:actualpolarizations}, where the subscript records
the prescribed polarization or its $4-t^2$ twist. Throughout this
subsection, Jacobian means the canonical principal polarization of a
smooth curve, up to origin-preserving polarized isomorphism.

\begin{proposition}[The balanced triple-cover count]
\label{pp:prop:trigonaltheta}
Let $C$ be a smooth genus-six cyclic triple cover of $\mathbb P^1$
with eight branch points, four of each nonzero inertia exponent.
If $\delta$ is a deck generator and $G=\ker(1-\delta)\subset J(C)$,
the unique $\delta$-invariant symmetric theta divisor has exactly
$213$ points of $G$ outside its support. Its multiplicities on $G$
are $0,1,2$ with respective frequencies $213,456,60$.
\end{proposition}
\begin{proof}
Move infinity away from the branch locus and write
$y^3=\prod_{i=1}^8(x-t_i)^{a_i}$, with
$a=(1,1,1,1,2,2,2,2)$. Let $P_i$ be the ramification points and
$H$ the fiber at infinity. Then $3P_i\sim H$ and
$D=\sum_iP_i-H$ has degree five and $2D\sim K_C$.
It is the unique invariant theta characteristic: the characteristic
torsor has translation group $J(C)[2]$, on which $1-\delta$ is
invertible since $(1-\delta)(2+\delta)=3$.

Put $W=\{b\in\mathbb F_3^8:\sum_i b_i=0\}$.
For representatives $b_i\in\{0,1,2\}$, define
\[
 E_b=\sum_i b_iP_i-\frac{\sum_i b_i}{3}H.
\]
The map $b\mapsto[E_b]$ takes values in $G$ and has kernel
$\langle a\rangle$. Indeed $a$ gives $\operatorname{div}(y)$;
conversely, a function with invariant principal divisor is a
$\delta$-eigenfunction $y^j h(x)$, and its orders at the $P_i$
force $b_i=ja_i$ modulo three. Both $W/\langle a\rangle$ and
$G$ have order $3^6$, the latter by the eigenvalues of $1-\delta$
on $H^1(C,\mathbb C)$. This identifies every point of $G$.

For $m=1+\sum_i b_i/3$, the three eigensheaves for sections of
$\mathcal O_C(D+E_b)$ on $\mathbb P^1$ have degrees
\[
 e_k=-m-4k+\sum_i\left\lfloor\frac{ka_i+1+b_i}{3}\right\rfloor,
 \qquad k=0,1,2.
\]
This follows by taking orders of $y^k h(x)$ at the branch points
and infinity. Thus $h^0(D+E_b)=\sum_k\max(e_k+1,0)$, independently
of the distinct branch positions. The exact finite count on the
$729$ classes gives sorted degree triples
\[
\begin{array}{c|rrrrr}
(e_0,e_1,e_2)_{\rm sorted}&(-1,-1,-1)&(-2,-1,0)&(-2,-2,1)&(-3,0,0)&(-3,-1,1)\\
\hline
\text{frequency}&213&456&24&24&12.
\end{array}
\]
The retained finite certificate is specified in
\path{experiments/bc_jacobian_v147/adoption.json}; its domain is
exactly the quotient just proved complete. Riemann's singularity
theorem identifies these dimensions with theta multiplicities
\cite[Introduction, smooth-curve statement]{CasalainaMartinKass2012}.
Translation by $D$ identifies the theta support in $\operatorname{Pic}^0(C)$
with $h^0(D+E_b)>0$. The principal line bundle has $h^0=1$,
so its invariant class determines a unique effective divisor.
The displayed frequencies give the assertion.
\end{proof}

\begin{theorem}[Jacobian classification under the enclosure contract]
\label{pp:thm:jacobiancomplete}
Assume the trusted Arb/acb enclosure and execution contract specified
below for the retained theta certificates. Exactly one of the ten
classes in Theorem~\ref{pp:thm:tenprincipal} is a canonically polarized
Jacobian of a smooth curve: the known class of the smooth projective
model of $y^7=x^3-1$. In particular, none of $B_0,B_1,C_0,C_1$ is
such a Jacobian.
\end{theorem}
\begin{proof}
For each model, $\delta=\zeta_3$ preserves the principal polarization
and satisfies $\delta^2+\delta+1=0$. Symmetric representatives form
a torsor under two-torsion, so the same invertibility argument gives
a unique invariant line-bundle class. A principal ample line bundle
has a one-dimensional space of sections, hence determines a unique
effective divisor in this class. Its support is stable under negation.
If a B/C model were a Jacobian, the necessary geometry in
Corollary~\ref{pp:cor:jacobianlist} would identify this divisor with
the one in Proposition~\ref{pp:prop:trigonaltheta}. Inverting the
deck generator does not change its fixed kernel or invariant divisor.

The retained certificates give, on each model separately, $107$
distinct nonzero classes modulo sign in $\ker(1-\delta)$ at which
the theta section is nonzero. The exact-coordinate sets are the
indexed representatives $0,\ldots,106$ in the retained pair map.
Their counts are
\[
\begin{array}{c|rrrr}
\text{model}&B_0&B_1&C_0&C_1\\\hline
\text{nonzero pairs}&107&107&107&107\\
\text{distinct points}&214&214&214&214.
\end{array}
\]
The kernel is killed by three, so a nonzero point differs from its
negative. Each model therefore supplies $214>213$ points outside
the invariant divisor, contradicting the Jacobian count. No origin
value or sum across different models is used. Combining these four
exclusions with Corollary~\ref{pp:cor:jacobianlist} proves the tally.
\end{proof}

The analytic computation starts from the exact ideal lattice, actual
CM embeddings and principal alternating form. It verifies an integral
symplectic basis, recomputes the $\delta$ action and its unique
invariant quadratic refinement in that basis, and constructs interval
periods. For normalized lattice $[I,\tau]$, a fixed point has the
form $z=(n+\tau m)/3$ with $(1-F)(n,m)^t=0$ modulo three.
The half-characteristic is read from that same refinement. These
conventions and their geometric identification are reviewed in
\path{experiments/bc_jacobian_v147/PERIOD_SETUP_REVIEW.md}.

The nonvanishing records are exact dyadic rectangles for squared theta
values. An independent rational check verifies that a real or imaginary
interval excludes zero and that all representatives are disjoint modulo
sign. This check does not reimplement the analytic evaluator: the
conclusion assumes correctness of the pinned Arb/acb routines, the
stated enclosure/serialization interface, and the admitted execution
and runtime provenance. An enclosure excluding zero for the square
implies nonvanishing of the section; it does not determine a zero's
multiplicity. The cumulative certificate is
\path{experiments/bc_jacobian_v147/certificate_run07b_adoption.json},
which binds the eight accepted retained packages; the geometric consequence
is bound by \path{experiments/bc_jacobian_v147/jacobian_classification_adoption.json}.
The separate hostile review in
\path{experiments/bc_jacobian_v147/session_v150/CLASSIFICATION_REVIEW03.md}
checks all $428$ raw rectangles and the complete implication above, including
the two explicitly admitted runtime variants. It does not prove the analytic
library implementation or its equivalence to source.

The theorem concerns the specified endomorphism-order category and
polarized isomorphism. It does not exclude isogenies to Jacobians or
address arbitrary CM orders, and establishes no additional Hodge
algebraicity or Hilbert-modular orbit membership.


% !TeX root = ../../quotients-and-descent.tex
\section{The smaller-order Jacobian cases}
The following criterion describes what a Jacobian in the $S$-order
family would have to satisfy. Partial nonvanishing computations do not
settle its $156$ polarized classes.

\begin{proposition}[A structural theta filter for balanced cubic actions]
\label{pp:prop:balancedtheta}
Let $(Y,\lambda)$ be a complex principally polarized abelian sixfold
with a polarized automorphism $\delta$ satisfying
$\delta^2+\delta+1=0$ and tangent signature $(3,3)$.
Its unique $\delta$-invariant symmetric principal line-bundle class
has a unique effective divisor $\Theta$. Normalize the order-three
bundle lift to fix the origin fiber. In $G=\ker(1-\delta)$ there
are $261$ points with trivial fiber character and $234$ with each
nontrivial character. The theta section has trivial character;
therefore the $468$ points in the nontrivial sectors lie on $\Theta$.
Every zero in the $261$-point sector is singular.
\end{proposition}
\begin{proof}
Write $\Lambda=H_1(Y,\Z)$, let $E$ be its unimodular Riemann form,
and let $F$ act as $\delta$. Since $(1-F)(2+F)=3$, the affine action
on the torsor of symmetric representatives has a unique fixed class.
Its semicharacter is $(-1)^{q(k)}$ for the invariant quadratic
refinement $q$ of $E$ modulo two. In the Appell--Humbert convention
$H(x,y)=E(x,Iy)-iE(x,y)$, where $I$ is the complex structure,
the multiplier is
$j_k(z)=(-1)^{q(k)}\exp(\pi H(z,k)+\pi H(k,k)/2)$;
see \cite[\S1.2, pp.~7--8]{PolishchukTheta}.
The lift $T(z,t)=(Fz,t)$ has order three and fixes the origin fiber.

Put $N=F-I_{12}$. Fixed points are
$x=-(F+2I_{12})k/3$, parametrized by $\Lambda/N\Lambda$.
This quotient has order $3^6$ and is killed by three, since
$N^2=-3F$ and $\det N=3^6$.
Set $b(k)=E(k,Fk)$. Invariance of $q$ gives
$q(k)=b(k)\pmod2$, and $E(3x,k)=b(k)$.
The normalized fiber character is consequently
\[
 j_k(x)^{-1}
 =\exp\bigl(\pi i(3q(k)+b(k))/3\bigr)
 =\zeta_3^{2b(k)}.
\]
Modulo three, the polar form of $b$ is
$B_b(k,l)=-E(k,Nl)$. Also $b(Nu)=0$ and $N^2=0$ modulo
three, so $b$ descends to $\Lambda/N\Lambda$.
Unimodularity makes the polar radical on $\Lambda/3\Lambda$
equal to $\ker N=\operatorname{im}N$, both of dimension six.
Thus the descended quadratic form is nondegenerate.

Its six-variable Gauss sum is $+27$ or $-27$: diagonalization
expresses it as a product of six sums $\pm i\sqrt3$.
The inverse lift $\delta^*L\to L$ has the conjugate fiber sum,
which has the same real value. Holomorphic Lefschetz
\cite[Theorem~4.12]{AtiyahBott1968II} applies, with simple fixed
points and complex denominator
$(1-\zeta_3)^3(1-\zeta_3^2)^3=27$.
Kodaira vanishing and Riemann--Roch give $H^{>0}(Y,L)=0$ and
$h^0(Y,L)=1$. The section eigenvalue is therefore $+1$ or $-1$;
its order divides three, so it must be $+1$. The Gauss sum is $27$.
Character orthogonality gives $261,234,234$ for the three levels
of $b$. Equivariance forces the nontrivial sectors to vanish.
At a zero in the trivial sector, the first derivative intertwines
the tangent action with the trivial fiber action. There is no
invariant tangent covector, so this derivative vanishes as well.
\end{proof}

\begin{corollary}[Necessary Jacobian geometry of the smaller-order targets]
\label{pp:cor:smallerjacobian}
If any of the $156$ $S$-order principal classes of
Theorem~\ref{pp:thm:fermatallquotients} is a canonical smooth-curve
Jacobian, its curve is nonhyperelliptic, is a cyclic triple cover of
$\mathbb P^1$ branched at eight points with inertia multiplicities
four and four, and has full automorphism group $C_3$.
For the invariant symmetric theta divisor, the fixed-point
multiplicity distribution is
\[
\begin{array}{c|rrr}
\text{fiber sector}&0&1&2\\\hline
\text{trivial }(261)&213&0&48\\
\text{nontrivial }(468)&0&456&12.
\end{array}
\]
The origin is outside the divisor. Among the $130$ eligible nonzero
opposite pairs there are then $106$ nonvanishing pairs and $24$
double-zero pairs.
\end{corollary}
\begin{proof}
Corollary~\ref{pp:cor:quotientweilhodge} supplies the integral cubic
action and signature, and Theorem~\ref{pp:thm:fermatallquotients}
gives the full polarized automorphism group $\mu_6$.
The signed Torelli, Riemann--Hurwitz, cyclic-cover eigenspace and
hyperelliptic arguments of Corollary~\ref{pp:cor:jacobianlist}
therefore apply with the same hypotheses, giving the asserted curve
geometry and full automorphism group.
Proposition~\ref{pp:prop:trigonaltheta} gives total multiplicities
$213,456,60$. Proposition~\ref{pp:prop:balancedtheta} puts all
nonvanishing points in the trivial sector and makes its zeros
singular. This gives the displayed table. At the origin, the three
eigensheaf degrees in the proof of
Proposition~\ref{pp:prop:trigonaltheta} are all $-1$, so its
invariant theta characteristic has no section. Negation now gives
the pair counts.
\end{proof}

The proof and regression are retained in
\path{experiments/s_order_jacobian_v160/}. In particular, the split
quadratic-form sign follows structurally from the balanced tangent
action; it is not extrapolated from the four B/C matrices.
No theta evaluation on these $156$ smaller-order classes is asserted,
and their Jacobian status remains open. A certificate of nonvanishing
on $107$ eligible nonzero opposite pairs would exclude the particular
principal object tested; the B/C certificates do not supply such
values on any new target.

The classical two-prime isogeny decomposition appears in
\cite[Eq.~(4.2) and Lem.~4.2]{Goodson2024}. The following argument
computes the restriction on the specified embedded primitive factor.


% !TeX root = ../../quotients-and-descent.tex
\section{Questions left by the classification}
Three finite questions remain particularly concrete.
\begin{enumerate}[leftmargin=*]
\item Which of the $156$ full-$S$ principal classes are canonical Jacobians?
The ten-class $R$-containing theorem does not answer this question.
A partial theta count is evidence only until the required threshold and
all model, enclosure and execution hypotheses have been met.
\item Are the definition-degree lower bounds $8$ and $4$ sharp for the
remaining $16$ and $4$ underlying classes, respectively? A minimum degree
is different from a unique smallest embedded field of definition.
\item For each of the $108$ obstructed varieties, what is the union of
principal-target kernels over its $26$ minimum polarization forms?
The fixed-form count cannot simply be multiplied by $26$, since kernels
may recur and target classes may coincide.
\end{enumerate}
These questions extend the proved classification and do not enter its
proofs. The construction supplies a finite testing ground for integral
CM geometry and descent; it asserts neither new points on McMullen's
curve nor a general classification of abelian sixfolds.

\begingroup\small
% !TeX root = ../quotients-and-descent.tex
\begin{thebibliography}{99}
\bibitem{AtiyahBott1968II}
M.~F.~Atiyah and R.~Bott, A Lefschetz fixed point formula for elliptic
complexes: II. Applications, \emph{Annals of Mathematics} (2)
\textbf{88} (1968), no.~3, 451--491; Theorem~4.12.
\url{https://webhomes.maths.ed.ac.uk/~v1ranick/papers/abott2.pdf}.

% PROVISIONAL: accompanying the separately reviewed ray-field claim.

\bibitem{CasalainaMartinKass2012}
S.~Casalaina-Martin and J.~L.~Kass,
\emph{A Riemann singularity theorem for integral curves},
Amer. J. Math. \textbf{134} (2012), no.~5, 1143--1165.
Introduction, classical smooth-curve statement; arXiv:0907.0212v2.
\url{https://arxiv.org/abs/0907.0212v2}

\bibitem{Goodson2024}
H. Goodson,
\emph{An exploration of degeneracy in abelian varieties of Fermat type},
arXiv:2211.03909v2, 1 August 2024.
The cited numbering refers to this version.
\url{https://arxiv.org/abs/2211.03909v2}

\bibitem{LiMantovanPries2021}
W.~Li, E.~Mantovan and R.~Pries,
\emph{Data for Shimura varieties intersecting the Torelli locus},
arXiv:2105.02286v1 (2021); Pacific J. Math. \textbf{343} (2026), 179--210.
Equation~(2.2), Proposition~6.11, Lemma~6.12 and Proposition~6.14
refer to the inspected arXiv version.
\url{https://arxiv.org/abs/2105.02286v1}

\bibitem{Masley1976}
J. M. Masley,
Solution of small class number problems for cyclotomic fields,
\emph{Compositio Mathematica} \textbf{33} (1976), 179--186,
\url{https://www.numdam.org/item/CM_1976__33_2_179_0.pdf}.

\bibitem{MilneModuli1972}
J. S. Milne,
Abelian varieties defined over their fields of moduli, I,
\emph{Bull. London Math. Soc.} \textbf{4} (1972), 370--372.
\url{https://www.jmilne.org/math/articles/1972c.pdf}

\bibitem{MilneModuli1974}
J. S. Milne,
Correction to ``Abelian varieties defined over their fields of moduli, I'',
\emph{Bull. London Math. Soc.} \textbf{6} (1974), 145--146,
including the addendum.

\bibitem{MilneAV2008}
J. S. Milne,
\emph{Abelian Varieties}, course notes, version 2.0, March 2008,
Chapter IV, Lemma 2.1, p.~135 (quotients by Galois-stable finite subgroups).
\url{https://www.jmilne.org/math/CourseNotes/AVc.pdf}

\bibitem{MilneShimura2017}
J.~S.~Milne, \emph{Introduction to Shimura Varieties},
revised September 16, 2017, Theorem~6.8 and Corollary~6.9,
\url{https://www.jmilne.org/math/xnotes/svi.pdf}.

\bibitem{MilneCFT2020}
J. S. Milne,
\emph{Class Field Theory}, version 4.03, 6 August 2020.
\url{https://www.jmilne.org/math/CourseNotes/CFT.pdf}

\bibitem{MilneCM2020}
J. S. Milne,
\emph{Complex Multiplication}, version 0.10, 14 July 2020,
Section 3.11 and Propositions 3.12--3.13 and 3.17, pp.~29--31.
\url{https://www.jmilne.org/math/CourseNotes/CM.pdf}

\bibitem{MilneJacobians2021}
J.~S.~Milne, \emph{Jacobian Varieties}, revised June12, 2021,
Theorem~12.1, p.~38; original chapter in \emph{Arithmetic Geometry}
(Springer, 1986), 167--212.
\url{https://www.jmilne.org/math/xnotes/JVs.pdf}

\bibitem{OpenAICM2026}
OpenAI math repository,
\emph{The rational Hodge conjecture for CM abelian varieties},
public manuscript dated 6 October 2026, introduction (main CM theorem)
and tensor-reduction section (auxiliary-factor contractions).
The comparison uses repository commit
\texttt{4eaee87548250b2682affda7d05d947c05554bae};
the general theorem and its formalization are not proof inputs here.
\url{https://github.com/openai/math/tree/4eaee87548250b2682affda7d05d947c05554bae/preprints/The-rational-Hodge-conjecture-for-CM-abelian-varieties-October-6-2026}

\bibitem{PolishchukTheta}
A.~Polishchuk, \emph{Abelian Varieties, Theta Functions and the Fourier
Transform}, Cambridge University Press, section~1.2, pp.~7--8.
Publisher excerpt: \url{https://assets.cambridge.org/97805218/08040/excerpt/9780521808040_excerpt.pdf}.

\bibitem{Saller2026}
K. Saller,
\emph{Split Weil structures on abelian sixfolds and the Hodge conjecture},
version 2, Zenodo record 23290280 (2026), Remarks~4.7--4.8.
\url{https://doi.org/10.5281/zenodo.23290280}

\bibitem{StacksDescent}
The Stacks Project Authors,
\emph{The Stacks Project}, Lemma 39.25.3 (tag 0CCJ),
and Section 35.35 (tag 02VV), finite locally free effective descent
and fpqc descent of morphisms.
\url{https://stacks.math.columbia.edu/tag/0CCJ};
\url{https://stacks.math.columbia.edu/tag/02VV}.

\bibitem{Zarhin2007}
Y.~G.~Zarhin, \emph{Cubic surfaces and cubic threefolds, jacobians
and intermediate jacobians}, arXiv:math/0610138v3 (2007), Theorem~1.1.
\url{https://math.nyu.edu/~tschinke/.manin/final/zarhin/zarhin.pdf}

\end{thebibliography}

\endgroup
\end{document}
